\SGAchapter{XII}{光滑仿射群概形的极大环面、Weyl 群、Cartan 子群及约化中心}{Tores maximaux, groupe de Weyl, sous-groupes de Cartan, centre réductif des schémas en groupes lisses et affines}
\SourcePage{141}
\begin{center}A. Grothendieck\end{center}
\footnote[0]{原本版本注（0）：2008 年 12 月 2 日 xy 版。}
从本讲起，与前面各讲不同，我们将使用\Term{algebraically-closed-field}{代数闭域} $k$ 上\Term{smooth}{光滑}\Term{algebraic-group}{代数群}结构的已知结果，尤其是 Chevalley 讨论班 1956/57 年《\Term{algebra}{代数} Lie 群的分类》中讲述的 Borel \Term{affine}{仿射}\Term{algebraic-group}{代数群}理论。\footnote[1]{编者注：此讨论班现已由 P. Cartier 修订，作为《Chevalley 全集》第 3 卷出版（Springer，2005 年）。}依照\Term{algebraic-group}{代数群}理论中的惯例，我们以 BIBLE 指称这一讨论班。后面几讲特别需要 BIBLE 第 4、5、6、7 讲的结果（BIBLE 后的数字表示讲次）。此外，\Term{scheme}{概形}理论似乎没有给 BIBLE 所述的 Borel 理论带来任何显著简化。因此没有必要在本讨论班中重新讲述这一理论；我们这里的目的是，从\Term{algebraically-closed-field}{代数闭域}上的已知结果\Term{pushout}{推出}任意基\Term{prescheme}{预概形}上成立的相应结果。（对于 Chevalley 的\Term{semisimple}{半单}\Term{group-structure}{群结构}理论，情形则不同；看来从任意基\Term{prescheme}{预概形}出发从头处理更为有利。）

在口头报告中（第 1 至 4 节沿用其内容），我们只考虑 $S$ 上的\Term{affine}{仿射}\Term{group-prescheme}{群预概形}，并且实质性地依赖第 XI 讲第 4 节的\Term{representability}{可表示性}\Term{theorem}{定理}。在本讲义中（见第 6 至 8 节），我们将简要说明，如何用一种更简单的方法消去\Term{affine}{仿射}性假设，而不使用第 XI 讲中最精细的结果。关于本讲结果的其他推广，也见第 XV、XVI 讲。（显然，第 XI、XII、XV、XVI 讲的内容应当彻底重新编排。）

\SGAsection{1}{极大环面}
\SourcePage{142}
\textbf{1.0。}先设 $G$ 是\Term{algebraically-closed-field}{代数闭域} $k$ 上的\Term{algebraic-group}{代数群}。称 $G$ 的一个\Term{algebraic-subgroup}{代数子群} $T$ 为\Term{maximal-torus}{极大环面}，如果它是\Term{torus}{环面}（这里因为 $k$ \Term{algebraically-closed}{代数闭}，这就是说它\Term{isomorphism}{同构}于形如 $\mathbf G_m^r$ 的群），且关于此性质是极大的。注意，由于 $k$ \Term{perfect}{完美}，$G_{\mathrm{réd}}$ 是 $G$ 的一个\Term{subgroup}{子群}，在 $k$ 上\Term{smooth}{光滑}，因而实质上是 BIBLE 意义下的\Term{algebraic-group}{代数群}。此外，$G$ 的任何\Term{reduced}{既约}\Term{subgroup}{子群}（例如\Term{torus}{环面}）自动是 $G_{\mathrm{réd}}$ 的\Term{subgroup}{子群}，所以 $G$ 的\Term{maximal-torus}{极大环面}恰好是 $G_{\mathrm{réd}}$ 的\Term{maximal-torus}{极大环面}。当 $G$ 是\Term{affine}{仿射}的，从而 $G_{\mathrm{réd}}$ 也是\Term{affine}{仿射}的时，Borel 的一个基本\Term{theorem}{定理}说，$G$ 的任意两个\Term{maximal-torus}{极大环面}都由 $G(k)=G_{\mathrm{réd}}(k)$ 中的元素\Term{conjugacy}{共轭}（BIBLE 6，\Term{theorem}{定理} 4 c)）；特别地，它们的\Term{dimension}{维数}相同。我们称 $G$ 的\Term{maximal-torus}{极大环面}的公共\Term{dimension}{维数}为 $G$ 的\Term{reductive-rank}{约化秩}。其实 $G$ 为\Term{affine}{仿射}的\Term{restriction}{限制}并不必要；这由 Chevalley 的一个已知\Term{theorem}{定理}\Term{pushout}{推出}，该\Term{theorem}{定理}断言，$k$ 上的任意\Term{connected}{连通}\Term{smooth}{光滑}\Term{algebraic-group}{代数群}都是\Term{abelian-variety}{阿贝尔簇}被\Term{affine}{仿射}群的\Term{extension}{扩张}；见第 6 节。第 1 至 4 节大多只讨论\Term{base}{基底}上的\Term{affine}{仿射}\Term{group-prescheme}{群预概形}。

设 $G$ 是 $k$ 上的\Term{smooth}{光滑}\Term{algebraic-group}{代数群}，$T$ 是 $G$ 的\Term{maximal-torus}{极大环面}；称 $T$ 在 $G$ 中的\Term{centralizer}{中心化子} $C(T)$ 为与 $T$ 对应的 $G$ 的\Term{cartan-subgroup}{Cartan 子群}。对我们而言，这是由 VIII 6.7 定义的\Term{group-subprescheme}{群子预概形}。但应注意，由于 $G$ 在 $k$ 上\Term{smooth}{光滑}，由 XI 5.3，\Term{cartan-subgroup}{Cartan 子群} $C$ 也如此；所以此时 $C$ 是 $G$ 中唯一在 $k$ 上\Term{smooth}{光滑}、且使 $C(k)$ 为 $G(k)$ 中 $T(k)$ 的\Term{centralizer}{中心化子}的\Term{group-subprescheme}{群子预概形}（即 BIBLE 意义下的\Term{algebraic-subgroup}{代数子群}）。也就是说，它实质上就是 BIBLE 意义下的\Term{centralizer}{中心化子}。
由上述\Term{conjugacy}{共轭}\Term{theorem}{定理}，与各个\Term{maximal-torus}{极大环面}对应的\Term{cartan-subgroup}{Cartan 子群}互相\Term{conjugacy}{共轭}，因而\Term{dimension}{维数}相同。我们称其公共\Term{dimension}{维数}为 $G$ 的\Term{nilpotent-rank}{幂零秩}；它等于 $G_{\mathrm{réd}}$ 的\Term{nilpotent-rank}{幂零秩}。以 $\rho_r(G)$、$\rho_n(G)$ 分别表示 $G$ 的\Term{reductive-rank}{约化秩}、\Term{nilpotent-rank}{幂零秩}，则
\[\rho_r(G)\leq\rho_n(G),\]
而其差
\[\rho_u(G)=\rho_n(G)-\rho_r(G)=\dim C/T\]
在 $G$ 为\Term{affine}{仿射}时可以称为 $G$ 的\Term{unipotent-rank}{幺幂秩}。当 $G$ \Term{smooth}{光滑}、\Term{affine}{仿射}且\Term{connected}{连通}时，$C$ 是\Term{nilpotent}{幂零}且\Term{connected}{连通}的\Term{algebraic-group}{代数群}（BIBLE 6，\Term{theorem}{定理} 6 a)、c)），故（BIBLE 6，\Term{theorem}{定理} 2）\Term{isomorphism}{同构}于乘积 $C_s\times C_u$，其中 $C_s=T$ 是起初的\Term{maximal-torus}{极大环面}（当然也是 $C$ 的\Term{maximal-torus}{极大环面}），而 $C_u$ 是\Term{smooth}{光滑}\Term{unipotent}{幺幂}\Term{subgroup}{子群}，即\Term{isomorphism}{同构}于 $\mathbf G_a$ 的群的逐次\Term{extension}{扩张}（BIBLE 6，\Term{theorem}{定理} 1 推论 1 及 BIBLE 7，\Term{theorem}{定理} 4）。所以此时也有
\[\rho_u(G)=\dim C_u.\]
\textbf{注记 1.1。}除了刚才对\Term{affine}{仿射}\Term{algebraic-group}{代数群}给出的三种秩，还有两种有用的秩，即\Term{semisimple-rank}{半单秩} $\rho_s(G)$，按定义它是\Term{quotient-group}{商群} $G/R$ 的\Term{reductive-rank}{约化秩}，其中 $R$ 为 $G$ 的\Term{radical}{根基}；以及\Term{infinitesimal-rank}{无穷小秩} $\rho_i(G)$，按定义它是 $G$ 的\Term{lie-algebra}{Lie 代数}的\Term{nilpotent-rank}{幂零秩}（将在下一讲定义并研究）。本讲不使用它们。
\SourcePage{143}
这里只指出不等式
\[\rho_s\leq\rho_r\leq\rho_n\leq\rho_i.\]
\begin{SGAresult}{\Term{lemma}{引理}}{1.2}
设 $G$ 是\Term{algebraically-closed-field}{代数闭域} $k$ 上的\Term{algebraic-group}{代数群}，$T$ 是 $G$ 的\Term{algebraic-subgroup}{代数子群}，$k'$ 是 $k$ 的\Term{algebraically-closed}{代数闭}\Term{field-extension}{扩域}，$G'$ 和 $T'$ 是由 $G$、$T$ \Term{base-change}{基变换}得到的群。则 $T$ 是 $G$ 的\Term{maximal-torus}{极大环面}，当且仅当 $T'$ 是 $G'$ 的\Term{maximal-torus}{极大环面}。
\end{SGAresult}
这是“\Term{finite-extension}{有限扩张}原理”EGA IV 9.1.1 的直接推论。
\begin{SGAresult}{定义}{1.3}
设 $S$ 是\Term{prescheme}{预概形}，$G$ 是\Term{finite-type}{有限型} $S$-\Term{group-prescheme}{群预概形}，$T$ 是 $G$ 的\Term{group-subprescheme}{群子预概形}。称 $T$ 为 $G$ 的\Term{maximal-torus}{极大环面}，如果：

 a) $T$ 是\Term{torus}{环面}（IX 1.3）；且

 b) 对每个 $s\in S$，以 $\bar s$ 表示 $\kappa(s)$ 的一个\Term{algebraic-closure}{代数闭包}的谱，则 $T_{\bar s}$ 是 $G_{\bar s}$ 中的\Term{maximal-torus}{极大环面}。
\end{SGAresult}
\textbf{注记 1.4。}由 1.2，当 $S$ 是\Term{algebraically-closed-field}{代数闭域}的谱时，这恢复了通常的定义；而且性质 1.3 对任意\Term{base-change}{基变换}\Term{stable}{稳定}。注意，由 X 8.8，在定义 1.3 中可将条件 a) 换成：

 a') $T$ 在 $S$ 上\Term{finite-presentation}{有限表示}且\Term{flat}{平坦}。

应当注意，定义 1.3 意义下的\Term{maximal-torus}{极大环面}确实在 $G$ 的所有\Term{subtorus}{子环面}中极大（这由 IX 2.9 立即\Term{pushout}{推出}）；但反过来，譬如在\Term{base}{基底} $S$ \Term{connected}{连通}时，只有当 $G$ 确实存在 1.3 意义下的\Term{maximal-torus}{极大环面}，此逆命题才可能正确，而一般并非如此（即使 $S=\operatorname{Spec}(\mathbf Z)$ 且 $G$“\Term{semisimple}{半单}”；另见 1.6）。不过我们将在第 XIV 讲看到，当 $S$ 为 Artin 的，或 $S$ 为\Term{local-scheme}{局部}\Term{scheme}{概形}且 $G$“\Term{reductive}{约化}”时，逆命题成立：此时 $G$ 的任意\Term{torus}{环面}都包含在某个\Term{maximal-torus}{极大环面}中。
\begin{SGAresult}{定义}{1.5}
设 $G$ 是域 $k$ 上的\Term{algebraic-group}{代数群}。称 $G_{\bar k}$ 的\Term{reductive-rank}{约化秩}（相应地，\Term{nilpotent-rank}{幂零秩}、\Term{unipotent-rank}{幺幂秩}等）为 $G$ 的相应秩，其中 $\bar k$ 为 $k$ 的一个\Term{algebraic-closure}{代数闭包}。
\end{SGAresult}
注意，由 1.2 以及 $\operatorname{Centr}_G(T)$ 的形成与\Term{base-change}{基变换}可交换，1.5 引入的各种秩在\Term{base-field-extension}{基域扩张}下\Term{invariant}{不变}；另一方面，当 $k$ \Term{algebraically-closed}{代数闭}时，它们与本节开头引入的秩一致。

\textbf{注记 1.6。}不难在一个\Term{discrete-valuation-ring}{离散赋值环}的谱 $S$ 上构造\Term{affine}{仿射}\Term{smooth}{光滑}\Term{group-scheme}{群概形} $G$，使其一般\Term{fiber}{纤维}\Term{isomorphism}{同构}于 $\mathbf G_m$，特殊\Term{fiber}{纤维}\Term{isomorphism}{同构}于 $\mathbf G_a$。这样的 $G$ 除了平凡\Term{torus}{环面} $T$（即单位\Term{subgroup}{子群}）以外不含任何\Term{torus}{环面}，而 $T$ 显然不是\Term{maximal-torus}{极大环面}。具体地，在特殊\Term{fiber}{纤维} $G_0=\mathbf G_{a,k}$ 中，$T_0$ 确为\Term{maximal-torus}{极大环面}；但在一般\Term{fiber}{纤维} $G_1=\mathbf G_{m,K}$ 中，$T_1$ 就不再是\Term{maximal-torus}{极大环面}（$k$ 是所考察\Term{valuation-ring}{赋值环}的\Term{residue-field}{剩余域}，$K$ 是其\Term{field-of-fractions}{分式域}）。
\SourcePage{144}
这个例子也表明，$G_s$ 的\Term{reductive-rank}{约化秩}（$s\in S$）不一定是 $s$ 的\Term{continuous}{连续}函数。不过有以下结果。
\begin{SGAresult}{\Term{theorem}{定理}}{1.7}
设 $S$ 为\Term{prescheme}{预概形}，$G$ 是 $S$ 上\Term{affine}{仿射}且\Term{smooth}{光滑}的 $S$-\Term{group-prescheme}{群预概形}。对每个 $s\in S$，以 $\rho_r(s)=\rho_r(G_s)$、$\rho_n(s)=\rho_n(G_s)$ 分别表示 $G_s$ 的\Term{reductive-rank}{约化秩}和\Term{nilpotent-rank}{幂零秩}（1.5）。在这些记号下：

 a) $S$ 上的函数 $\rho_r$ \Term{lower-semicontinuous}{下半连续}，函数 $\rho_n$ \Term{upper-semicontinuous}{上半连续}，因而函数 $\rho_u=\rho_n-\rho_r$ \Term{upper-semicontinuous}{上半连续}。

 b) 以下条件\Term{equivalence}{等价}，且对任意\Term{base-change}{基变换}\Term{stable}{稳定}：

 (i) $S$ 上的函数 $\rho_r$ \Term{locally-constant}{局部常值}。

 (ii) 按\Term{etale-topology}{平展拓扑}\Term{local-scheme}{局部}地，$G$ 存在\Term{maximal-torus}{极大环面}。

 (ii bis) 按\Term{fpqc-topology}{忠实平坦拟紧拓扑}\Term{local-scheme}{局部}地，$G$ 存在\Term{maximal-torus}{极大环面}。

 c) 设 $T_1$、$T_2$ 是 $G$ 中两个\Term{maximal-torus}{极大环面}（这意味着 b) 中的条件成立）。则它们按\Term{etale-topology}{平展拓扑}\Term{local-scheme}{局部}\Term{conjugacy}{共轭}；即存在\Term{surjectivity}{满射}\Term{etale}{平展}\Term{morphism}{态射} $S'\to S$，使 $G_{S'}$ 的\Term{subgroup}{子群} $(T_1)_{S'}$、$(T_2)_{S'}$ 由 $G_{S'}$ 在 $S'$ 上的一个\Term{section}{截面}\Term{conjugacy}{共轭}。

 d) 在 b) 的条件下，即当 $\rho_r$ \Term{locally-constant}{局部常值}时，$\rho_n$ 也\Term{locally-constant}{局部常值}，因而 $\rho_u=\rho_n-\rho_r$ 也如此。
\end{SGAresult}
\begin{proof}
a) 注意，对任意\Term{morphism}{态射} $S'\to S$，若 $G'=G\times_S S'$，则按 $G'$ 定义的 $S'$ 上的函数 $\rho'_r$ 等，只需将由 $G$ 定义的函数 $\rho_r$ 等与 $S'\to S$ 复合即可得到。当 $S'\to S$ \Term{faithfully-flat}{忠实平坦}且\Term{quasi-compact}{拟紧}时，因 $S$ 的 Zariski \Term{topology}{拓扑}是 $S'$ 的 Zariski \Term{topology}{拓扑}的商（SGA 1 VIII 4.3），$\rho'$ \Term{continuous}{连续}、\Term{upper-semicontinuous}{上半连续}、\Term{lower-semicontinuous}{下半连续}分别当且仅当 $\rho$ 具有相应性质。因此 a) 中的断言对于\Term{fpqc-topology}{忠实平坦拟紧拓扑}是\Term{local-scheme}{局部}的。
设 $s\in S$，要证明由满足 $\rho_r(t)\geq\rho_r(s)$（相应地，$\rho_n(t)\leq\rho_n(s)$）的 $t\in S$ 组成的\Term{set}{集合} $U$ 是 $s$ 的\Term{neighborhood}{邻域}。由\Term{finite-extension}{有限扩张}原理，存在 $\kappa(s)$ 的\Term{finite}{有限}\Term{field-extension}{扩域} $k$，使 $G_k$ 存在\Term{maximal-torus}{极大环面}。于是存在 $s$ 的\Term{open-neighborhood}{开邻域} $U$ 和\Term{finite}{有限}、\Term{surjectivity}{满射}、\Term{locally-free}{局部自由}的\Term{morphism}{态射} $S'\to U$，使\Term{fiber}{纤维} $S'_s$ 在 $\kappa(s)$ 上\Term{isomorphism}{同构}于 $\operatorname{Spec}(k)$（见 EGA III 10.3.2，其中 Noether 假设显然不必要）。由于 $S'\to S$ 是\Term{open-morphism}{开态射}（SGA 1 IV 6.6），可归结到 $S'=S$ 的情形，即 $G_s$ 存在\Term{maximal-torus}{极大环面} $T_s$。此外，由 XI 5.8 a)，再以一个在 $S$ 上\Term{etale}{平展}且带有 $s$ 上方一点 $s'$ 的 $S'$ 替换 $S$ 后，可设 $T_s$ 是 $G$ 的某个\Term{subtorus}{子环面} $T$ 在 $s$ 处的\Term{fiber}{纤维}。于是对每个 $t\in S$，
\[\rho_r(t)=\rho_r(G_t)\geq\dim T_t=\dim T_s=\rho_r(G_s)=\rho_r(s),\]
证明了 $\rho_r$ \Term{lower-semicontinuous}{下半连续}。另一方面，由 XI 5.3，\Term{functor}{函子}
\[\mathcal C=\operatorname{Centr}_G(T)\]
\SourcePage{145}
由 $C=\operatorname{Centr}_G(T)$ \Term{representable}{表示}，\footnote[2]{编者注：这里作了修改，引入表示\Term{functor}{函子} $\operatorname{Centr}_G(T)$ 的\Term{prescheme}{预概形} $C=\operatorname{Centr}_G(T)$。}它是 $G$ 的\Term{closed}{闭}\Term{group-subprescheme}{群子预概形}，且在 $S$ 上\Term{smooth}{光滑}。所以存在 $s$ 的\Term{neighborhood}{邻域} $U$，使 $t\in U$ 蕴含 $\dim C_t=\dim C_s=\rho_n(s)$。$\rho_n$ 的\Term{upper-semicontinuity}{上半连续性}就由关系
\[\rho_n(t)\leq\dim C_t\qquad(t\in S)\]
\Term{pushout}{推出}；而此关系又包含在下面这个纯几何\Term{lemma}{引理}中。
\end{proof}
\begin{SGAresult}{\Term{lemma}{引理}}{1.8}
设 $k$ 为域，$G$ 是 $k$ 上的\Term{affine}{仿射}\Term{smooth}{光滑}\Term{algebraic-group}{代数群}，$T$ 是 $G$ 中的\Term{torus}{环面}，$C$ 是其\Term{centralizer}{中心化子}。则
\[\rho_n(G)\leq\dim C.\]
\end{SGAresult}
\begin{proof}
可设 $k$ \Term{algebraically-closed}{代数闭}，于是 $T$ 包含在某个\Term{maximal-torus}{极大环面} $T'$ 中。设 $C'$ 是后者的\Term{centralizer}{中心化子}，则 $C'\subset C$，因而 $\dim C'\leq\dim C$。
\end{proof}
\begin{proof}[\Term{theorem}{定理} 1.7 的证明（续）]
b) 若 $\rho_r$ \Term{locally-constant}{局部常值}，则对 $G$ 中的任意\Term{torus}{环面} $T$ 及任意 $s\in S$，若 $T_s$ 是 $G_s$ 中的\Term{maximal-torus}{极大环面}，就存在 $s$ 的\Term{open-neighborhood}{开邻域} $U$，使 $T|_U$ 是 $G|_U$ 中的\Term{maximal-torus}{极大环面}。现在应用 a) 的论证即得 (i)$\Rightarrow$(ii bis)。另一方面，(ii bis)$\Rightarrow$(i)，因为当 $G$ 存在\Term{maximal-torus}{极大环面} $T$ 时，$\rho_r(s)=\dim T_s$ 显然是 $s$ 的\Term{locally-constant}{局部常值}函数，而在 a) 证明开头已指出，$\rho_r$ 的\Term{continuity}{连续性}对于\Term{faithfully-flat}{忠实平坦}\Term{topology}{拓扑}是\Term{local-scheme}{局部}问题。因此 (i)$\Leftrightarrow$(ii bis)，而显然 (ii)$\Rightarrow$(ii bis)，只须证明逆向蕴含 (i)$\Rightarrow$(ii)。为此引入 XI 4.1 的\Term{functor}{函子} $F$，\footnote[3]{编者注：这是 $G$ 的\Term{multiplicative-type}{乘法型}\Term{subgroup}{子群}\Term{functor}{函子}。}它是 $S$ 上\Term{smooth}{光滑}且\Term{separated}{分离}的\Term{prescheme}{预概形}；考察 $F$ 的\Term{subfunctor}{子函子} $\mathcal T$，其在 $S$ 上任意 $S'$ 处的值为 $G_{S'}$ 中\Term{maximal-torus}{极大环面}的\Term{set}{集合}。利用前面的观察：在 (i) 下，若 $G_{S'}$ 的一个\Term{torus}{环面}在一点 $s\in S'$ 的\Term{fiber}{纤维}中极大，则它在 $s$ 的某个\Term{open-neighborhood}{开邻域}上也极大，可见 $\mathcal T$ 由 $F$ 的一个开\Term{subprescheme}{子预概形}\Term{representable}{表示}，从而在 $S$ 上\Term{smooth}{光滑}且\Term{separated}{分离}。\Term{structural-morphism}{结构态射} $\mathcal T\to S$ 显然\Term{surjectivity}{满射}，所以由 XI 1.10，按\Term{etale-topology}{平展拓扑}\Term{local-scheme}{局部}地存在\Term{section}{截面}，这就证明了 (i)$\Rightarrow$(ii)。

 c) 这是 XI 5.4 bis 的直接推论，只须考虑本节开头回顾的 Borel \Term{conjugacy}{共轭}\Term{theorem}{定理}。

 d) 由 a) 证明开头的观察及 b)，可设 $G$ 存在\Term{maximal-torus}{极大环面} $T$。若 $C$ 是其\Term{centralizer}{中心化子}，则由 XI 5.3，$C$ \Term{representable}{可表示}且在 $S$ 上\Term{smooth}{光滑}，所以函数 $s\mapsto\rho_n(s)=\dim C_s$ 确为\Term{locally-constant}{局部常值}。

\Term{theorem}{定理} 1.7 得证。今后称满足 1.7 b) 中条件的 $G$ 具有\Term{locally-constant}{局部常值}\Term{reductive-rank}{约化秩}。记下：
\end{proof}
\SourcePage{146}
\begin{SGAresult}{推论}{1.9}
设 $G$ 如 1.7，且 $s\in S$ 满足 $\rho_u(s)=0$，即 $\rho_r(s)=\rho_n(s)$（也就是说，$G_k$ 的\Term{cartan-subgroup}{Cartan 子群}都是\Term{torus}{环面}，其中 $k$ 是 $\kappa(s)$ 的一个\Term{algebraic-closure}{代数闭包}）。则 $s$ 有一个\Term{open-neighborhood}{开邻域} $U$，使 $\rho_r$、$\rho_n$ 在其上为常值；特别地，每个 $t\in U$ 处，$G_t$ 的\Term{unipotent-rank}{幺幂秩} $\rho_u(t)$ 都为零。
\end{SGAresult}
事实上，这由 1.7 a) 和对所有 $t\in S$ 成立的不等式 $\rho_r(t)\leq\rho_n(t)$ 立即\Term{pushout}{推出}。还应指出，在证明 b) 的同时，我们已经证明了：
\begin{SGAresult}{推论}{1.10}
设 $G$ 如 1.7，且 $G$ 的\Term{reductive-rank}{约化秩}\Term{locally-constant}{局部常值}。考虑\Term{functor}{函子}
\[\mathcal T:(\mathrm{Sch})^\circ\longrightarrow(\mathrm{Ens}),\qquad
\mathcal T(S')=\{G_{S'}\text{ 的极大环面}\}.\]
则 $\mathcal T$ 由 $S$ 上\Term{smooth}{光滑}、\Term{separated}{分离}且\Term{finite-type}{有限型}的\Term{prescheme}{预概形}\Term{representable}{表示}。
\end{SGAresult}
\begin{proof}
只须证明 $\mathcal T$ 在 $S$ 上\Term{finite-type}{有限型}。此问题对于\Term{fpqc-topology}{忠实平坦拟紧拓扑}是\Term{local-scheme}{局部}的，所以可设 $G$ 存在\Term{maximal-torus}{极大环面} $T$。由 XI 5.3 bis 和 5.5 的 bis 版本，$\operatorname{Norm}_G(T)$ 及 $G/\operatorname{Norm}_G(T)$ 由\Term{prescheme}{预概形} $\operatorname{Norm}_G(T)$ 及 $G/\operatorname{Norm}_G(T)$ \Term{representable}{表示}，且 $\mathcal T$ \Term{isomorphism}{同构}于 $G/\operatorname{Norm}_G(T)$。\footnote[4]{编者注：这里作了修改，引入\Term{prescheme}{预概形} $\operatorname{Norm}_G(T)$ 和 $G/\operatorname{Norm}_G(T)$。}由 $g\mapsto\operatorname{int}(g)(T)$ 定义的\Term{morphism}{态射} $G\to\mathcal T$ \Term{surjectivity}{满射}，而 $G$ 在 $S$ 上\Term{quasi-compact}{拟紧}，所以 $\mathcal T$ 也如此，证明完毕。
\end{proof}
\textbf{注记 1.11。}a) 1.10 的\Term{prescheme}{预概形} $\mathcal T$ 自然称为 $G$ 的\Term{maximal-torus-prescheme}{极大环面预概形}。我们将在第 5 节看到，它事实上在 $S$ 上\Term{affine}{仿射}。当 $G$ \Term{isomorphism}{同构}于某个形如 $\mathrm{GL}(n)$ 的\Term{prescheme}{预概形}的\Term{closed}{闭}\Term{group-subprescheme}{群子预概形}时，可直接用 XI 4.6 证明这一点：由 1.7 b) 证明中的论证，$\mathcal T$ 等同于表示 $G$ 的\Term{multiplicative-type}{乘法型}\Term{subgroup}{子群}的\Term{prescheme}{预概形} $F$ 的一个既开又闭的\Term{subprescheme}{子预概形}。

 b) 可以给出 1.7（从而 1.9）的证明，只用 XI 3.12 和 6.2 的简单结果，不用第 XI 讲中的精细结果；办法是只处理\Term{base}{基底}上\Term{finite}{有限}的\Term{multiplicative-type}{乘法型}群（大体上即 $T$ 为\Term{maximal-torus}{极大环面}时的群 ${}_nT$），比较第 7 节。对 1.10 的证明也有同样的说明。

最后给出存在唯一\Term{maximal-torus}{极大环面}的例子，以结束本节。
\begin{SGAresult}{命题}{1.12}
设 $G$ 是\Term{multiplicative-type}{乘法型}且\Term{finite-type}{有限型}的 $S$-\Term{group-prescheme}{群预概形}。则 $G$ 存在唯一\Term{maximal-torus}{极大环面}，而 $G$ 中每个\Term{torus}{环面}都包含在这一\Term{maximal-torus}{极大环面}中。
\end{SGAresult}
\SourcePage{147}
\begin{proof}
唯一性显然由后一断言\Term{pushout}{推出}；后一断言将此\Term{maximal-torus}{极大环面}刻画为 $G$ 的最大\Term{subtorus}{子环面}。由唯一性，存在性对于\Term{fpqc-topology}{忠实平坦拟紧拓扑}是\Term{local-scheme}{局部}问题，因此可设 $G$ \Term{diagonalizable}{可对角化}，即形如 $\mathrm D_S(M)$，其中 $M$ 是\Term{finite-type}{有限生成}\Term{abelian-group}{交换群}。设 $M_0$ 是 $M$ 除以其\Term{torsion-subgroup}{挠子群}的商。我断言，$G$ 中的\Term{torus}{环面} $T=\mathrm D_S(M_0)$ 是\Term{maximal-torus}{极大环面}，也是最大\Term{subtorus}{子环面}。事实上，$G$ 的\Term{subtorus}{子环面} $T'$ 按 \Term{fpqc-topology}{fpqc 拓扑}\Term{local-scheme}{局部}\Term{diagonalizable}{可对角化}；故为证明 $T'\subset T$，可设 $T'$ \Term{diagonalizable}{可对角化}，于是形如 $\mathrm D_S(N)$，其中 $N$ 为 $M$ 的自由商（VIII 1.4 和 3.2 b)），所以 $N$ 是 $M_0$ 的商，故 $T'\subset T$。因为将 $T$ 构造为 $\mathrm D_S(M_0)$ 与任意\Term{base-change}{基变换}相容，这同时说明 $T$ 是 $G$ 的\Term{maximal-torus}{极大环面}，证明完毕。当 $G$ 在 $S$ 上\Term{smooth}{光滑}时，可将 1.12 推广如下。
\end{proof}
\begin{SGAresult}{命题}{1.13}
设 $G$ 是在 $S$ 上\Term{finite-presentation}{有限表示}的 $S$-\Term{group-prescheme}{群预概形}。假定按\Term{fpqc-topology}{忠实平坦拟紧拓扑}\Term{local-scheme}{局部}地，$G$ 存在一个\Term{central}{中心}的\Term{maximal-torus}{极大环面}。则 $G$（整体上）存在唯一\Term{maximal-torus}{极大环面}，\footnote[5]{编者注：而且它是\Term{central}{中心}的！}并且它是 $G$ 的最大\Term{subtorus}{子环面}。
\end{SGAresult}
唯一性同样显然由后一断言\Term{pushout}{推出}，并使存在性成为 fpqc \Term{local-scheme}{局部}的问题。因此可设 $G$ 有\Term{central}{中心}\Term{maximal-torus}{极大环面} $T$。要证明 $G$ 的任意\Term{torus}{环面} $R$ 包含在 $T$ 中，这由下面\Term{lemma}{引理}\Term{pushout}{推出}。
\begin{SGAresult}{\Term{lemma}{引理}}{1.14}
设 $G$ 是在 $S$ 上\Term{finite-presentation}{有限表示}的 $S$-\Term{group-prescheme}{群预概形}，$T$ 是 $G$ 的\Term{maximal-torus}{极大环面}，$R$ 是 $G$ 的\Term{subtorus}{子环面}。若 $T$ 与 $R$ 交换，则 $R\subset T$。
\end{SGAresult}
\begin{proof}
因为 $R$ 与 $T$ 交换，由 $(r,t)\mapsto r\cdot t$ 定义的\Term{morphism}{态射} $R\times_S T\to G$ 是\Term{group-homomorphism}{群同态}，所以由 IX 6.8，它在 $G$ 中有像\Term{subgroup}{子群} $T'$，是 $R\times_S T$ 的\Term{multiplicative-type}{乘法型}\Term{quotient-group}{商群}，因而是\Term{torus}{环面}，且显然包含 $T$。由于 $T$ 是\Term{maximal-torus}{极大环面}，有 $T'=T$（1.4），故 $R\subset T$。这证明了 1.14，从而证明了 1.13。特别地，由 1.7 b) 得：
\end{proof}
\begin{SGAresult}{推论}{1.15}
设 $G$ 是 $S$ 上交换、\Term{smooth}{光滑}、\Term{affine}{仿射}且\Term{reductive-rank}{约化秩}\Term{locally-constant}{局部常值}的 $S$-\Term{group-prescheme}{群预概形}。则 $G$ 存在唯一\Term{maximal-torus}{极大环面}，且它包含 $G$ 的每个\Term{subtorus}{子环面}。
\end{SGAresult}
\begin{SGAresult}{推论}{1.16}
设 $G$ 是 $S$ 上\Term{smooth}{光滑}且\Term{affine}{仿射}的 $S$-\Term{group-prescheme}{群预概形}。假定对每个 $s\in S$，以 $\bar s$ 表示 $\kappa(s)$ 的一个\Term{algebraic-closure}{代数闭包} $k$ 的谱，\Term{geometric-fiber}{几何纤维} $G_{\bar s}$ 都是\Term{connected}{连通}、\Term{nilpotent}{幂零}的\Term{algebraic-group}{代数群}（BIBLE 意义下，即其 $k$ 值点群 $G(k)$ \Term{nilpotent}{幂零}）。再假定 $G$ 的\Term{reductive-rank}{约化秩}\Term{locally-constant}{局部常值}。则 $G$ 存在唯一\Term{maximal-torus}{极大环面} $T$，且 $T$ 是\Term{central}{中心}的，也是 $G$ 的最大\Term{subtorus}{子环面}。
\end{SGAresult}
\SourcePage{148}
\begin{proof}
由 1.7 b)，$G$ 按 fpqc \Term{local-scheme}{局部}地存在\Term{maximal-torus}{极大环面}；由 1.13，只须证明 $G$ 的\Term{maximal-torus}{极大环面}是\Term{central}{中心}的。由 IX 5.6 b)，可归结到 $S$ 是域的谱的情形，并可设此域\Term{algebraically-closed}{代数闭}。于是由 BIBLE 6 \Term{theorem}{定理} 2，$T(k)$ 包含在 $G(k)$ 的\Term{center}{中心}中。这蕴含 $T$ 包含在 $G$ 的\Term{center}{中心}中，因为 $\operatorname{Centr}_G(T)$ 是 $G$ 的\Term{closed}{闭}\Term{subscheme}{子概形}（VIII 6.6），包含 $G(k)$ 的点，从而由 Hilbert \Term{nullstellensatz}{零点定理}等于 $G$（因为 $G$ \Term{reduced}{既约}）。
\end{proof}
最后，为供后面引用，指出下列平凡命题（我们已暗中使用过它）。
\begin{SGAresult}{命题}{1.17}
设 $G\supset H\supset T$ 是 $S$ 上\Term{finite-type}{有限型}\Term{group-prescheme}{群预概形}。下列条件\Term{equivalence}{等价}：

 (i) $T$ 是 $G$ 的\Term{maximal-torus}{极大环面}。

 (ii) $T$ 是 $H$ 的\Term{maximal-torus}{极大环面}，且对每个 $s\in S$，\Term{reductive-rank}{约化秩}相等：$\rho_r(G_s)=\rho_r(H_s)$。
\end{SGAresult}
\SGAsection{2}{Weyl 群}
先设 $k$ 是\Term{algebraically-closed-field}{代数闭域}，$G$ 是 $k$ 上\Term{smooth}{光滑}且\Term{affine}{仿射}的\Term{algebraic-group}{代数群}。若 $T$ 是\Term{maximal-torus}{极大环面}，$C$、$N$ 分别是它的\Term{centralizer}{中心化子}、\Term{normalizer}{正规化子}，则由 XI 5.9，二者都是 $G$ 的\Term{smooth}{光滑}\Term{closed}{闭}\Term{subgroup}{子群}，且 $C$ 是 $N$ 的\Term{open-subgroup}{开子群}；所以 $W=N/C$ 是 $k$ 上的\Term{finite}{有限}\Term{etale}{平展}群，因而由它的 $k$ 值点群 $W(k)$ 决定，具体为 $W=W(k)_k$（由通常意义下的\Term{finite-group}{有限群} $W(k)$ 定义的\Term{constant-group}{常值群}）。\Term{finite-group}{有限群} $W(k)$ 称为 $G$ 相对于 $T$ 的\Term{geometric-weyl-group}{几何 Weyl 群}（若不致混淆，也简称\Term{weyl-group}{Weyl 群}）。由 Borel \Term{conjugacy}{共轭}\Term{theorem}{定理}，各个\Term{maximal-torus}{极大环面}的\Term{weyl-group}{Weyl 群}彼此\Term{isomorphism}{同构}，因此有时不指定\Term{maximal-torus}{极大环面}，而称为 $G$ 的“那个”\Term{weyl-group}{Weyl 群}。由于 $C$、$N$、$N/C$ 的形成与任意\Term{base-extension}{基底扩张}可交换，若 $k'$ 是 $k$ 的\Term{algebraically-closed}{代数闭}\Term{field-extension}{扩域}，则 $G_{k'}$ 相对于 $T_{k'}$ 的\Term{geometric-weyl-group}{几何 Weyl 群}与 $G$ 相对于 $T$ 的相应群\Term{canonical-isomorphism}{典范同构}；所以 $G$ 的“那个”\Term{geometric-weyl-group}{几何 Weyl 群}（严格说，是通常\Term{finite-group}{有限群}的一个\Term{isomorphism-class}{同构类}）与 $G'$ 的相同。

由此，当 $G$ 是任意域 $k$ 上\Term{smooth}{光滑}且\Term{affine}{仿射}的\Term{algebraic-group}{代数群}时，可将 $G$ 的\Term{geometric-weyl-group}{几何 Weyl 群}定义为 $G_{k'}$ 的\Term{geometric-weyl-group}{几何 Weyl 群}的\Term{isomorphism-class}{同构类}，其中 $k'$ 是 $k$ 的任意\Term{algebraically-closed}{代数闭}\Term{field-extension}{扩域}。当 $G$ 有\Term{maximal-torus}{极大环面} $T$ 时，显然可如上构造 $C=\operatorname{Centr}_G(T)$、$N=\operatorname{Norm}_G(T)$、$W=N/C$；这是 $k$ 上的\Term{finite}{有限}\Term{etale}{平展}群，称为 $G$ 相对于 $T$ 的\Term{weyl-group}{Weyl 群}。此时\Term{geometric-weyl-group}{几何 Weyl 群}就是 $W$ 在任意\Term{algebraically-closed}{代数闭}\Term{field-extension}{扩域} $k'$ 中的点群的\Term{isomorphism-class}{同构类}。这里，通常仅知道\Term{geometric-weyl-group}{几何 Weyl 群} $W(k')$ 显然已不足以恢复\Term{algebraic-group}{代数群} $W$：还必须知道 $k$ 在 $k'$ 中的\Term{separable}{可分}\Term{algebraic-closure}{代数闭包}的 Galois 群
\SourcePage{149}
在 $W(k')$ 上的\Term{action}{作用}。

最后，若 $G$ 是任意\Term{base}{基底} $S$ 上的\Term{group-prescheme}{群预概形}，在 $S$ 上\Term{smooth}{光滑}且\Term{affine}{仿射}，而 $T$ 是 $G$ 的\Term{maximal-torus}{极大环面}，则 XI 5.9 仍容许构造群
\[W(T)=N(T)/C(T),\]
它是 $S$ 上\Term{etale}{平展}、\Term{separated}{分离}、\Term{quasi-finite}{拟有限}的 $S$-\Term{group-prescheme}{群预概形}。它的\Term{geometric-fiber}{几何纤维}（相对于\Term{residue-field}{剩余域} $\kappa(s)$ 的\Term{algebraic-closure}{代数闭包}，$s\in S$）就是各\Term{fiber}{纤维} $G_s$ 的\Term{geometric-weyl-group}{几何 Weyl 群}。因此 XI 5.10 给出了这些群随 $s\in S$ 变化的一些信息。可以如下加以明确和推广。
\begin{SGAresult}{\Term{theorem}{定理}}{2.1}
设 $S$ 是\Term{prescheme}{预概形}，$G$ 是 $S$ 上\Term{smooth}{光滑}且\Term{affine}{仿射}的 $S$-\Term{group-prescheme}{群预概形}。对每个 $s\in S$，设 $w(s)$ 是 $G_s$ 的\Term{geometric-weyl-group}{几何 Weyl 群}，即\Term{finite-group}{有限群}的一个\Term{isomorphism-class}{同构类}。在这些\Term{isomorphism-class}{同构类}的\Term{set}{集合} $E$ 上引入\Term{preorder}{预序}：$w\leq w'$ 当且仅当 $w$、$w'$ 分别由\Term{finite-group}{有限群} $W$、$W'$ 代表，且 $W$ \Term{isomorphism}{同构}于 $W'$ 的某个\Term{subgroup}{子群}的商。在此约定下：

 a) 函数 $s\mapsto w(s)$，$S\to E$，\Term{lower-semicontinuous}{下半连续}。

 b) 假定 $G$ 的\Term{reductive-rank}{约化秩}\Term{locally-constant}{局部常值}。则以下条件\Term{equivalence}{等价}：

 (i) 函数 $s\mapsto w(s)$ \Term{locally-constant}{局部常值}。

 (i bis) 函数 $s\mapsto\operatorname{card}w(s)$ \Term{locally-constant}{局部常值}。

 (ii) 按\Term{etale-topology}{平展拓扑}（或仅按\Term{fpqc-topology}{fpqc 拓扑}）\Term{local-scheme}{局部}地，存在\Term{maximal-torus}{极大环面} $T$，使 $W(T)$ 在其\Term{base}{基底}上\Term{finite}{有限}。

 (ii bis) 对 $S$ 上每个 $S'$ 和 $G_{S'}$ 的每个\Term{maximal-torus}{极大环面} $T$，对应的\Term{weyl-group}{Weyl 群} $W(T)$ 在 $S'$ 上\Term{finite}{有限}。
\end{SGAresult}
\begin{proof}
a) 如 1.7 a) 那样，为证明对每个 $s\in S$ 存在 $s$ 的\Term{open-neighborhood}{开邻域} $U$，使 $t\in U$ 蕴含 $w(t)\geq w(s)$，可归结到 $G$ 中存在\Term{torus}{环面} $R$，使 $R_s$ 是 $G_s$ 中\Term{maximal-torus}{极大环面}的情形。如 XI 5.9，置
\[W(R)=\operatorname{Norm}_G(R)/\operatorname{Centr}_G(R),\]
它是 $S$ 上\Term{etale}{平展}、\Term{separated}{分离}、\Term{quasi-finite}{拟有限}的\Term{group-prescheme}{群预概形}。对每个 $t\in S$，以 $w'(t)\in E$ 表示它在 $t$ 处的\Term{geometric-fiber}{几何纤维}。因为 $R_s$ 是 $G_s$ 中的\Term{maximal-torus}{极大环面}，而 $\operatorname{Norm}$、$\operatorname{Centr}$、$\operatorname{Norm}/\operatorname{Centr}$ 的形成与任何\Term{base-extension}{基底扩张}、特别是取\Term{fiber}{纤维}相容，故有
\[w(s)=w'(s).\]
我还断言，对邻近 $s$ 的 $t$，有
\[w(t)\geq w'(t)\quad\text{和}\quad w'(t)\geq w'(s),\]
这足以证明 a)。这两个不等式包含在以下两个\Term{lemma}{引理}中。
\end{proof}
\SourcePage{150}
\begin{SGAresult}{\Term{lemma}{引理}}{2.2}
设 $S$ 是\Term{prescheme}{预概形}，$W$ 是 $S$ 上\Term{etale}{平展}、\Term{separated}{分离}、\Term{quasi-finite}{拟有限}的 $S$-\Term{group-prescheme}{群预概形}。对每个 $s\in S$，设 $f(s)$ 是 $W$ 在 $s$ 处的\Term{geometric-fiber}{几何纤维}的\Term{isomorphism-class}{同构类}，视为 2.1 引入的\Term{finite-group}{有限群}\Term{isomorphism-class}{同构类}的\Term{ordered-set}{有序集} $E$ 中的元素。则函数 $f:S\to E$ \Term{lower-semicontinuous}{下半连续}。它在 $s\in S$ 的某\Term{neighborhood}{邻域}上为常值，当且仅当函数 $s\mapsto\operatorname{card}f(s)$ 也如此；又当且仅当存在 $s$ 的\Term{open-neighborhood}{开邻域} $U$，使 $W$ 在 $U$ 上\Term{finite}{有限}。
\end{SGAresult}
\begin{proof}
这个结果以群的语言精细化了 XI 5.10 证明中所引的结果。我们只略述证明（它属于最标准的类型）。照常可归结到 $S$ \Term{affine}{仿射}且 Noether 的情形。立即看出函数 $f$ \Term{constructible}{可构造}（EGA $0_{\mathrm{III}}$ 9.3.1、9.3.2 及 EGA IV 第 9 节的基本性质）；由 EGA $0_{\mathrm{III}}$ 9.3.4，要证明\Term{semicontinuity}{半连续性}，只须证明当 $t$ 是 $s$ 的\Term{generization}{泛化}时，有 $f(t)\geq f(s)$。借助 EGA II 7.1.7，可归结到 $S$ 为\Term{discrete-valuation-ring}{离散赋值环}的谱；还可设此环\Term{complete}{完备}且\Term{residue-field}{剩余域}\Term{algebraically-closed}{代数闭}。但此时以 $s$ 表示 $S$ 的\Term{closed-point}{闭点}，因为 $G$ 在 $S$ 上\Term{etale}{平展}且\Term{separated}{分离}，它包含一个既开又闭、在 $S$ 上\Term{finite}{有限}的\Term{subscheme}{子概形} $G'$，满足 $G'_s=G_s$（EGA III 6.2.6），并且立即看出，这里的 $G'$ 是 $G$ 的\Term{subgroup}{子群}。此外，$G'$ 在 $S=\operatorname{Spec}(V)$ 上\Term{finite}{有限}\Term{etale}{平展}，而 $V$ \Term{complete}{完备}且\Term{residue-field}{剩余域}\Term{algebraically-closed}{代数闭}，故 $G'$ 为\Term{constant-group}{常值群}，形如 $A_S$，其中 $A=G'(\kappa(s))=G(\kappa(s))$ 的\Term{isomorphism-class}{同构类}为 $f(s)$。若 $B$ 为 $G$ 在 $S$ 的\Term{generic-point}{一般点} $t$ 处的\Term{geometric-fiber}{几何纤维}，就有\Term{canonical}{典范}\Term{monomorphism}{单态射} $A\to B$，这证明 $f(s)\leq f(t)$。（注意，此证明实际上对 $E$ 上比 2.1 所述更细的\Term{order-relation}{序关系}也证明了\Term{semicontinuity}{半连续性}。）$f$ 在 $s$ 处\Term{continuous}{连续}当且仅当 $s\mapsto\operatorname{card}f(s)$ \Term{continuous}{连续}，这是因为对 $w,w'\in E$，关系 $w\leq w'$ 及 $\operatorname{card}w=\operatorname{card}w'$ 蕴含 $w=w'$。此条件\Term{equivalence}{等价}于 $W$ 在 $s$ 的某\Term{neighborhood}{邻域}上\Term{finite}{有限}，这一事实与 $W$ 的\Term{group-structure}{群结构}无关，且已在 XI 5.10 之后指出；也容易用前述论证加上\Term{proper}{固有}性的赋值判别法 EGA II 7.3.8 来证明。
\end{proof}
\begin{SGAresult}{\Term{lemma}{引理}}{2.3}
设 $G$ 是\Term{algebraically-closed-field}{代数闭域} $k$ 上的\Term{affine}{仿射}\Term{smooth}{光滑}\Term{algebraic-group}{代数群}，$R\subset T$ 是两个\Term{subtorus}{子环面}，$W(R)$、$W(T)$ 是如 XI 5.9 所定义的相应\Term{finite-group}{有限群}，即\Term{normalizer}{正规化子}除以\Term{centralizer}{中心化子}的商。则 $W(R)$ \Term{isomorphism}{同构}于 $W(T)$ 的某个\Term{subgroup}{子群}的商。
\end{SGAresult}
\begin{proof}
考虑图表
\[\begin{tikzcd}
T\arrow[r,hook]&C(T)\arrow[r,hook]\arrow[d,hook]&C(R)\arrow[d,hook]\\
&N(T)\cap N(R)\arrow[r,hook]\arrow[d,hook]&N(R)\\
&N(T)&
\end{tikzcd}\]
\SourcePage{151}
于是 $(N(T)\cap N(R))/C(T)$ 是 $W(T)=N(T)/C(T)$ 的\Term{subgroup}{子群}，并且有显然的\Term{homomorphism}{同态}
\[(N(T)\cap N(R))/C(T)\longrightarrow W(R)=N(R)/C(R).\]
只须证明它\Term{surjectivity}{满射}，即对 $N(R)$ 的每个 $k$ 值点 $g$，存在 $C(R)$ 的一个 $k$ 值点 $c$，使 $cg$ \Term{normalization}{正规化} $T$，也就是
\[\operatorname{int}(c)(\operatorname{int}(g)T)=T.\]
为此，只须注意 $\operatorname{int}(g)T$ 是 $N(R)$ 中的\Term{torus}{环面}，因而也是 $C(R)$ 中的\Term{torus}{环面}（后者是前者的\Term{open-subgroup}{开子群}）。于是 $T$ 和 $\operatorname{int}(g)T$ 是 $C(R)$ 的\Term{maximal-torus}{极大环面}，因为它们在 $G$ 中极大；再用 Borel \Term{conjugacy}{共轭}\Term{theorem}{定理}即可。
\end{proof}
这证明了 2.3，从而证明了 2.1 a)。
\begin{proof}[\Term{theorem}{定理} 2.1 的证明（续）]
b) 我们已指出，(i) 和 (i bis) 显然\Term{equivalence}{等价}；由 2.2 的逆向结论或 XI 5.10，它们蕴含 (ii bis)；由 1.7 b)，(ii bis)$\Rightarrow$(ii)。最后，(ii)$\Rightarrow$(i bis)：如 1.7 a) 一样，条件 (i) 对于 fpqc 是\Term{local-scheme}{局部}的，所以可设 $G$ 存在\Term{maximal-torus}{极大环面} $T$，且 $W(T)$ 在 $S$ 上\Term{finite}{有限}；再次用 XI 5.10 即得结论。这就完成了 2.1 的证明。
\end{proof}
\SGAsection{3}{Cartan 子群}
\begin{SGAresult}{定义}{3.1}
设 $G$ 是\Term{prescheme}{预概形} $S$ 上\Term{smooth}{光滑}\Term{finite-type}{有限型}的\Term{group-prescheme}{群预概形}。称 $G$ 的\Term{group-subprescheme}{群子预概形} $C$ 为 $G$ 的\Term{cartan-subgroup}{Cartan 子群}，如果 $C$ 在 $S$ 上\Term{smooth}{光滑}，且对每个 $s\in S$，以 $\bar s$ 表示 $\kappa(s)$ 的一个\Term{algebraic-closure}{代数闭包}的谱，$C_{\bar s}$ 是 $G_{\bar s}$ 的\Term{cartan-subgroup}{Cartan 子群}（见第 1 节）。
\end{SGAresult}
显然，若 $C$ 是 $G$ 的\Term{cartan-subgroup}{Cartan 子群}，则对 $S$ 上每个 $S'$，$C_{S'}$ 是 $G_{S'}$ 的\Term{cartan-subgroup}{Cartan 子群}。若 $S$ 是\Term{algebraically-closed-field}{代数闭域}的谱，就恢复了第 1 节回顾的概念。最后立即可见，$G$ 的\Term{group-subprescheme}{群子预概形} $C$ 为\Term{cartan-subgroup}{Cartan 子群}，是关于\Term{fpqc-topology}{忠实平坦拟紧拓扑}的\Term{local-nature}{局部性质}。
\begin{SGAresult}{\Term{theorem}{定理}}{3.2}
设 $G/S$ 如 3.1，再假定 $G$ 在 $S$ 上\Term{affine}{仿射}且\Term{reductive-rank}{约化秩}\Term{locally-constant}{局部常值}。则映射
\[T\longmapsto\operatorname{Centr}_G(T)\]
给出 $G$ 的\Term{maximal-torus}{极大环面}\Term{set}{集合}与 $G$ 的\Term{cartan-subgroup}{Cartan 子群}\Term{set}{集合}之间的\Term{bijection}{双射}。\footnote{原注（*）：这里给出的证明实际上只对 $G$ 的\Term{closed}{闭}\Term{cartan-subgroup}{Cartan 子群}证明了 3.2。不过 7.1 a) 按所述\Term{form}{形式}建立了 3.2，并蕴含 $G$ 的\Term{cartan-subgroup}{Cartan 子群}都是\Term{closed}{闭}的。}若 $C$ 对应于 $T$，则 $T$ 是 $C$ 的唯一\Term{maximal-torus}{极大环面}。
\end{SGAresult}
\SourcePage{152}
\begin{proof}
若 $T$ 是 $G$ 的\Term{maximal-torus}{极大环面}，则由 XI 5.3 或 6.2，\Term{functor}{函子} $\operatorname{Centr}_G(T)$ 由 $G$ 的\Term{closed}{闭}\Term{smooth}{光滑}\Term{subprescheme}{子预概形} $C=\operatorname{Centr}_G(T)$ \Term{representable}{表示}；由定义，$C$ 是 $G$ 的\Term{cartan-subgroup}{Cartan 子群}。此外，$T$ 显然是 $C$ 的\Term{maximal-torus}{极大环面}，而因为它是\Term{central}{中心}的，便是 $C$ 的唯一\Term{maximal-torus}{极大环面}（1.13）。所以映射 $T\mapsto\operatorname{Centr}_G(T)$ 从\Term{maximal-torus}{极大环面}\Term{set}{集合}到\Term{cartan-subgroup}{Cartan 子群}\Term{set}{集合}是\Term{injectivity}{单射}，只须证明\Term{surjectivity}{满射}。于是设 $C$ 是 $G$ 的\Term{cartan-subgroup}{Cartan 子群}，要证明 $C=\operatorname{Centr}_G(T)$，其中 $T$ 是 $G$ 的\Term{maximal-torus}{极大环面}。为此只须找到 $C$ 的\Term{maximal-torus}{极大环面} $T$，因为它也是 $G$ 的\Term{maximal-torus}{极大环面}（对每个 $s\in S$，$G_s$ 和 $C_s$ 的\Term{reductive-rank}{约化秩}相同）。由 IX 5.6 b)，$T$ 在 $C$ 的\Term{center}{中心}内，故 $C\subset C'=\operatorname{Centr}_G(T)$；于是 $C$ 是 $S'$ 上\Term{smooth}{光滑}群 $C'$ 的\Term{smooth}{光滑}\Term{subgroup}{子群}，逐\Term{fiber}{纤维}与 $C'$ 相同，所以 $C=C'$。又因 $G$ 的\Term{reductive-rank}{约化秩}\Term{locally-constant}{局部常值}，$C$ 也如此，故由 1.7 b)，$C$ 按\Term{etale-topology}{平展拓扑}\Term{local-scheme}{局部}地存在\Term{maximal-torus}{极大环面}。由前面的论证，此\Term{torus}{环面}是\Term{central}{中心}的；故由 1.13，$C$ 确实存在\Term{maximal-torus}{极大环面}，证明完毕。
\end{proof}
\begin{SGAresult}{推论}{3.3}
设 $G$ 是 $S$ 上\Term{smooth}{光滑}、\Term{affine}{仿射}且\Term{reductive-rank}{约化秩}\Term{locally-constant}{局部常值}的\Term{group-prescheme}{群预概形}。令 $\mathcal C:(\mathrm{Sch})^\circ_{/S}\to(\mathrm{Ens})$ 为\Term{functor}{函子}
\[\mathcal C(S')=\{G_{S'}\text{ 的 Cartan 子群}\}.\]
则\Term{functor}{函子} $\mathcal C$ \Term{isomorphism}{同构}于 1.10 的\Term{functor}{函子} $\mathcal T$，所以由同一个在 $S$ 上\Term{smooth}{光滑}、\Term{separated}{分离}、\Term{finite-type}{有限型}的\Term{prescheme}{预概形}\Term{representable}{表示}。
\end{SGAresult}
\begin{SGAresult}{推论}{3.4}
在 3.2 的条件下，若 $C=\operatorname{Centr}_G(T)$，则
\[\operatorname{Norm}_G(C)=\operatorname{Norm}_G(T).\]
\end{SGAresult}
\textbf{注记 3.5。}当 $G$ 的\Term{reductive-rank}{约化秩}并非\Term{locally-constant}{局部常值}时，$G$ 仍可能存在\Term{cartan-subgroup}{Cartan 子群}（例如若 $G$ 的\Term{fiber}{纤维}\Term{connected}{连通}且\Term{nilpotent}{幂零}，则 $G$ 是自身的\Term{cartan-subgroup}{Cartan 子群}，但其\Term{reductive-rank}{约化秩}未必\Term{locally-constant}{局部常值}，见 1.6）。在第 XV 讲，我们将利用下一讲的\Term{regular-element}{正则元素}理论，建立\Term{cartan-subgroup}{Cartan 子群}理论，而不假定 $G$ 在 $S$ 上\Term{affine}{仿射}，也不假定其\Term{reductive-rank}{约化秩}\Term{locally-constant}{局部常值}。关于消去\Term{affine}{仿射}性假设，另见第 6、7 节。
\SGAsection{4}{约化中心}
\begin{SGAresult}{定义}{4.1}
设 $S$ 是\Term{prescheme}{预概形}，$G$ 是在 $S$ 上\Term{finite-presentation}{有限表示}且\Term{fiber}{纤维}\Term{affine}{仿射}的 $S$-群，$Z$ 是 $G$ 的\Term{group-subprescheme}{群子预概形}。称 $Z$ 为 $G$ 的\Term{reductive-center}{约化中心}，如果：

 (i) $Z$ 是\Term{central}{中心}且\Term{multiplicative-type}{乘法型}的；

 (ii) 对每个\Term{base-change}{基变换} $S'\to S$ 及每个\Term{central}{中心}\Term{homomorphism}{同态} $u:H\to G_{S'}$，其中 $H$ 为 $S$ 上\Term{multiplicative-type}{乘法型}\Term{finite-type}{有限型}群，$u$ 都经由 $Z_{S'}$ \Term{resolution}{分解}。

（不假定 $G$ 的\Term{fiber}{纤维}\Term{affine}{仿射}时，这一概念的一种变体见 8.6。）
\end{SGAresult}
\SourcePage{153}
注意，这样的 $Z$ 必在 $S$ 上\Term{finite-type}{有限型}（因为其\Term{fiber}{纤维}如此），所以作为 $G$ 的最大\Term{central}{中心}\Term{multiplicative-type}{乘法型}\Term{subgroup}{子群}而唯一确定。容易举出例子，使 $G$（在 $S$ 上\Term{smooth}{光滑}且\Term{affine}{仿射}）存在最大\Term{central}{中心}\Term{multiplicative-type}{乘法型}\Term{subgroup}{子群} $Z$，但 $Z$ 并非\Term{reductive-center}{约化中心}（即 $G$ 不存在\Term{reductive-center}{约化中心}）；例如见 1.6。不过由 IX 6.8 容易\Term{pushout}{推出}，$G$ 的\Term{subgroup}{子群} $Z$ 是 $G$ 的\Term{reductive-center}{约化中心}，当且仅当它是最大\Term{central}{中心}\Term{multiplicative-type}{乘法型}\Term{subgroup}{子群}，且在任意\Term{base-change}{基变换}下保持此性质；另见下面的 4.3。

由 4.1 显然，若 $Z$ 是 $G$ 的\Term{reductive-center}{约化中心}，则对每个\Term{base-change}{基变换} $S'\to S$，$Z_{S'}$ 是 $G_{S'}$ 的\Term{reductive-center}{约化中心}。由此及\Term{reductive-center}{约化中心}的唯一性，借助\Term{faithfully-flat}{忠实平坦}\Term{quasi-compact}{拟紧}\Term{descent}{下降}理论（SGA 1 VIII），可得：
\begin{SGAresult}{命题}{4.2}
设 $G$ 是 $S$ 上\Term{finite-presentation}{有限表示}且\Term{fiber}{纤维}\Term{affine}{仿射}的 $S$-\Term{group-prescheme}{群预概形}。若 $G$ 按\Term{fpqc-topology}{忠实平坦拟紧拓扑}\Term{local-scheme}{局部}地存在\Term{reductive-center}{约化中心}，则 $G$ 存在\Term{reductive-center}{约化中心}。$Z$ 是 $G$ 的\Term{reductive-center}{约化中心}，当且仅当按\Term{fpqc-topology}{fpqc 拓扑}\Term{local-scheme}{局部}地如此。
\end{SGAresult}
还记下：
\begin{SGAresult}{命题}{4.3}
设 $G$ 是 $S$ 上\Term{finite-presentation}{有限表示}且\Term{fiber}{纤维}\Term{affine}{仿射}的 $S$-\Term{group-prescheme}{群预概形}，$Z$ 是其\Term{group-subprescheme}{群子预概形}。则 $Z$ 是 $G$ 的\Term{reductive-center}{约化中心}，当且仅当 $Z$ 为\Term{multiplicative-type}{乘法型}，且对每个 $s\in S$，$Z_s$ 是 $G_s$ 的\Term{reductive-center}{约化中心}。
\end{SGAresult}
\begin{proof}
首先由 IX 5.6 b)，此时 $Z$ 是\Term{central}{中心}的。由于所述条件对\Term{base-change}{基变换}\Term{stable}{稳定}，只须证明任意\Term{central}{中心}\Term{homomorphism}{同态} $u:H\to G$（$H$ 为\Term{multiplicative-type}{乘法型}且\Term{finite-type}{有限型}）经由 $Z$ \Term{resolution}{分解}。因为 $Z$ 是\Term{central}{中心}的，$u$ 与\Term{canonical}{典范}\Term{immersion}{浸入} $v:Z\to G$ 定义\Term{group-homomorphism}{群同态}
\[w:H\times_S Z\longrightarrow G.\]
由 IX 6.8，此\Term{homomorphism}{同态}有像群 $K$，它是 $G$ 的\Term{multiplicative-type}{乘法型}\Term{subgroup}{子群}，只须证明 $K=Z$。由对 $Z$ 的假设，逐\Term{fiber}{纤维}确实如此，故应用 IX 5.1 bis 即完成 4.3 的证明。
\end{proof}
注意，在判别准则 4.3 中，对每个 $s\in S$，$Z_s$ 是 $G_s$ 的\Term{reductive-center}{约化中心}这一假设实际上纯属几何条件；即只须在 $\kappa(s)$ 的\Term{algebraic-closure}{代数闭包}上核查，这由 4.2 的第二个断言\Term{pushout}{推出}。
\begin{SGAresult}{\Term{theorem}{定理}}{4.4}
设 $G$ 是域 $k$ 上的\Term{affine}{仿射}\Term{algebraic-group}{代数群}。则 $G$ 存在\Term{reductive-center}{约化中心}。
\end{SGAresult}
\begin{proof}
因为 $G$ 的\Term{center}{中心}由 $G$ 的\Term{closed}{闭}\Term{subgroup}{子群}\Term{representable}{表示}（VIII 6.7），立即可归结到 $G$ 交换的情形。此外由 4.2，可设 $k$ \Term{algebraically-closed}{代数闭}。我们将在第 XVII 讲看到，此时 $G$ 可写为乘积 $Z\times U$，其中 $Z$ 为\Term{multiplicative-type}{乘法型}，而 $U$“\Term{unipotent}{幺幂}”；
\SourcePage{154}
这就立即说明 $Z$ 确为 $G$ 的\Term{reductive-center}{约化中心}。这里给出一个独立于上述一般\Term{form}{形式}的交换\Term{algebraic-group}{代数群}结构\Term{theorem}{定理}的证明。

由 IX 6.8，$G$ 的\Term{multiplicative-type}{乘法型}\Term{subgroup}{子群} $H$ 组成一个递增的\Term{filtered-family}{滤过族}。证明它有最大元。

当 $G$ 在 $k$ 上\Term{smooth}{光滑}时，应用结构\Term{theorem}{定理} BIBLE 4 \Term{theorem}{定理} 4，
\[G=G_s\times G_u,\]
其中 $G_s$ 为\Term{multiplicative-type}{乘法型}，$G_u$“\Term{unipotent}{幺幂}”，这里的意思是 $G_u$ 有一个\Term{composition-series}{合成列}，其因子为 $\mathbf G_a$ 的\Term{subgroup}{子群}（若愿意，还可以取\Term{smooth}{光滑}\Term{subgroup}{子群}）。（事实上，由 BIBLE 4 \Term{theorem}{定理} 3 的推论，$G_u/G_u^0$ 是\Term{unipotent}{幺幂}群，故由 BIBLE 4 命题 4，它是 $p$-群，其中 $p$ 为\Term{characteristic-exponent}{特征指数}。另一方面，由 BIBLE 6 \Term{theorem}{定理} 1 推论 1，$G_u^0$ 有一个\Term{composition-series}{合成列}，其商\Term{connected}{连通}、\Term{smooth}{光滑}且\Term{dimension}{维数}为 1；由 BIBLE 7 \Term{theorem}{定理} 4，这些商\Term{isomorphism}{同构}于 $\mathbf G_a$。）容易看出（见下面\Term{lemma}{引理}），从\Term{multiplicative-type}{乘法型}群到 $\mathbf G_a$，从而到 $G_u$ 的任何\Term{homomorphism}{同态}都是平凡的，故 $G$ 的任意\Term{multiplicative-type}{乘法型}\Term{subgroup}{子群}都包含在 $G_s$ 中。

在一般情形下，考虑 $G$ 的\Term{subgroup}{子群}
\[G_0=G_{\mathrm{réd}}.\]
它在 $k$ 上\Term{smooth}{光滑}，故由上面结论，它有最大\Term{multiplicative-type}{乘法型}\Term{subgroup}{子群} $Z_0$。$G$ 中包含 $Z_0$ 的\Term{multiplicative-type}{乘法型}\Term{subgroup}{子群}对应于 $G'=G/Z_0$ 的\Term{multiplicative-type}{乘法型}\Term{subgroup}{子群}。又有\Term{exact-sequence}{正合列}
\[0\longrightarrow G_0/Z_0\longrightarrow G/Z_0\longrightarrow G/G_0\longrightarrow0,\]
由前面的结论，$G_0/Z_0$ 除\Term{unit-group}{单位群}外没有\Term{multiplicative-type}{乘法型}\Term{subgroup}{子群}。因为\Term{multiplicative-type}{乘法型}群的\Term{subgroup}{子群}仍为\Term{multiplicative-type}{乘法型}（IX.8），所以对 $G/Z_0$ 的任意\Term{multiplicative-type}{乘法型}\Term{subgroup}{子群} $H$，有 $H\cap(G_0/Z_0)=0$，故 $H\to G/G_0$ 为\Term{injectivity}{单射}。因 $G/G_0$ 是\Term{radicial}{根性}\Term{algebraic-group}{代数群}，从而在 $k$ 上\Term{finite}{有限}，这意味着 $H$ 自身\Term{radicial}{根性}，且秩不超过 $G/G_0$ 的秩。因此 $G$ 的\Term{multiplicative-type}{乘法型}\Term{subgroup}{子群}族有最大元，记为 $Z$。

我断言，$G/Z$ 除单位\Term{subgroup}{子群}外没有\Term{multiplicative-type}{乘法型}\Term{subgroup}{子群}。这由以下事实（在 7.1.1 证明）\Term{pushout}{推出}：两个\Term{multiplicative-type}{乘法型}\Term{algebraic-group}{代数群}的交换\Term{extension}{扩张}仍为\Term{multiplicative-type}{乘法型}。因此，若 $Z'$ 是 $G/Z$ 的一个\Term{multiplicative-type}{乘法型}\Term{subgroup}{子群}在 $G$ 中的\Term{inverse-image}{逆像}，则 $Z'$ 为\Term{multiplicative-type}{乘法型}，故由 $Z$ 的极大性，$Z'=Z$，从而 $Z'/Z=0$。

现在由“\Term{finite-extension}{有限扩张}原理”容易\Term{pushout}{推出}，对 $k$ 的任何\Term{field-extension}{扩域} $K$，$(G/Z)_K=G_K/Z_K$ 除\Term{unit-group}{单位群}外也没有\Term{multiplicative-type}{乘法型}\Term{subgroup}{子群}。

于是可以证明 $Z$ 是 $G$ 的\Term{reductive-center}{约化中心}。设 $u:H\to G_S$ 是 $S$-\Term{group-homomorphism}{群同态}，其中 $S$ 是 $k$ 上的\Term{prescheme}{预概形}，$H$ 是 $S$ 上\Term{multiplicative-type}{乘法型}\Term{finite-type}{有限型}群。要证明 $u$ 经由 $Z_S$ \Term{resolution}{分解}；\Term{equivalence}{等价}地，复合\Term{homomorphism}{同态} $H\to G_S\to(G/Z)_S=G_S/Z_S$ 为零。
\SourcePage{155}
置 $U=G/Z$，我断言任意\Term{homomorphism}{同态} $u:H\to U_S$ 都为零。事实上，由 IX 5.2，可归结到 $S$ 是域的谱的情形；再由 IX 6.8，此结论由 $U_K$ 除 $0$ 外没有\Term{multiplicative-type}{乘法型}\Term{subgroup}{子群}这一事实\Term{pushout}{推出}。这就完成了 4.4 的证明。只须补上下面\Term{lemma}{引理}的证明。
\end{proof}
\begin{SGAresult}{\Term{lemma}{引理}}{4.4.1}
设 $H$ 是\Term{multiplicative-type}{乘法型}的 $S$-\Term{group-prescheme}{群预概形}，则任意 $S$-\Term{group-homomorphism}{群同态} $u:H\to\mathbf G_a$ 都是平凡的。
\end{SGAresult}
\begin{proof}
将模 $\mathcal O_S^2=E$ 视为 $\mathcal O_S=E'$ 被 $\mathcal O_S=E''$ 的\Term{extension}{扩张}，则 $\mathbf G_a$ 等同于这个\Term{extension}{扩张}的\Term{automorphism}{自同构}\Term{prescheme}{预概形}，所以\Term{homomorphism}{同态} $u:H\to\mathbf G_a$ 等同于 $E$ 上保持\Term{extension}{扩张}结构的 $H$-模结构，即 $E'$ 在 $H$ 下\Term{stable}{稳定}，且 $H$ 在 $E'$、$E''$ 上的\Term{induced-action}{诱导作用}平凡。由 I 4.7.3，$H$ 在 $E$ 上的\Term{action}{作用}也平凡，故 $u$ 平凡。
\end{proof}
\textbf{注记 4.4.2。}若 $G$ 是 $k$ 上非零\Term{abelian-variety}{阿贝尔簇}，则若在 4.1 的定义中省去“$G$ \Term{affine}{仿射}”这一\Term{restriction}{限制}，$G$ 并不存在该意义下的\Term{reductive-center}{约化中心}。因为当 $n$ 与\Term{characteristic}{特征}互素时，${}_nG$ 在 $k$ 上\Term{etale}{平展}，阶与\Term{characteristic}{特征}互素，所以为\Term{multiplicative-type}{乘法型}；而 ${}_nG$ 这一族在 $G$ 中\Term{schematically-dense}{概形稠密}，故若存在\Term{reductive-center}{约化中心}，它就应等于 $G$，但 $G$ 并非\Term{multiplicative-type}{乘法型}，矛盾。这就是 4.1 中应\Term{restriction}{限制} $G$ 的\Term{fiber}{纤维}为\Term{affine}{仿射}的原因。
\begin{SGAresult}{\Term{lemma}{引理}}{4.5}
设 $S$ 是\Term{prescheme}{预概形}，$G$ 是 $S$ 上\Term{smooth}{光滑}、\Term{affine}{仿射}且\Term{fiber}{纤维}\Term{connected}{连通}的 $S$-\Term{group-prescheme}{群预概形}，$T$ 是 $G$ 的\Term{maximal-torus}{极大环面}，$u:H\to G$ 是\Term{central}{中心}\Term{homomorphism}{同态}，其中 $H$ 为\Term{multiplicative-type}{乘法型}\Term{finite-type}{有限型} $S$-\Term{group-prescheme}{群预概形}。则 $u$ 经由 $T$ \Term{resolution}{分解}。
\end{SGAresult}
\begin{proof}
设 $C$ 为 $T$ 的\Term{centralizer}{中心化子}，它是 $G$ 的\Term{closed}{闭}\Term{group-subprescheme}{群子预概形}，在 $S$ 上\Term{smooth}{光滑}（XI 5.3），所以在 $S$ 上\Term{affine}{仿射}。由于 $T$ 在 $G$ 的\Term{center}{中心}内，它是\Term{normal-subgroup}{正规}的，故可考虑\Term{quotient-group}{商群} $G/T=U$，它是\Term{representable}{可表示}的（VIII 5.1）。因为 $u$ 为\Term{central}{中心}\Term{homomorphism}{同态}，它经由 $C$ \Term{resolution}{分解}；只须证明复合\Term{homomorphism}{同态} $H\to C\to U=C/T$ 平凡。由 IX 5.2，可归结到 $S$ 为域的谱的情形，并可设该域\Term{algebraically-closed}{代数闭}。此时由 BIBLE 6 \Term{theorem}{定理} 2，$U$ 是\Term{connected}{连通}（在 $k$ 上\Term{smooth}{光滑}）的“\Term{unipotent}{幺幂}”\Term{algebraic-group}{代数群}；如前所述，即有一个商\Term{isomorphism}{同构}于 $\mathbf G_a$ 的\Term{composition-series}{合成列}。因此从\Term{multiplicative-type}{乘法型}\Term{finite-type}{有限型}群 $H$ 到 $U$ 的任意\Term{homomorphism}{同态}都平凡。这证明了 4.5。
\end{proof}
\begin{SGAresult}{推论}{4.6}
设 $G$ 如 4.5。若 $G$ 存在\Term{reductive-center}{约化中心}，则此\Term{center}{中心}包含在 $G$ 的每个\Term{maximal-torus}{极大环面}中。
\end{SGAresult}
\begin{SGAresult}{\Term{theorem}{定理}}{4.7}
设 $S$ 是\Term{prescheme}{预概形}，$G$ 是 $S$ 上\Term{smooth}{光滑}、\Term{affine}{仿射}且\Term{fiber}{纤维}\Term{connected}{连通}的 $S$-\Term{group-prescheme}{群预概形}。

 a) 对每个 $s\in S$，设 $z(s)$ 是 $G_s$ 的\Term{reductive-center}{约化中心}的型（由 4.4，此\Term{center}{中心}有定义）。将\Term{finite-type}{有限生成} $\mathbf Z$-模\Term{isomorphism-class}{同构类}的\Term{set}{集合} $E$ 排序，规定 $M$ 的\Term{isomorphism-class}{同构类}大于 $M'$ 的\Term{isomorphism-class}{同构类}，当且仅当 $M'$ \Term{isomorphism}{同构}于 $M$ 的商。
\SourcePage{156}
则映射 $S\to E$，$s\mapsto z(s)$，\Term{lower-semicontinuous}{下半连续}。

 b) $G$ 存在\Term{reductive-center}{约化中心} $Z$，当且仅当前述函数 $z:S\to E$ \Term{locally-constant}{局部常值}。若如此，则 $G/Z$ \Term{representable}{可表示}（见 VIII 5.1），且 $G/Z$ 以单位\Term{subgroup}{子群}为\Term{reductive-center}{约化中心}。

 c) 假定 $G$ 的\Term{reductive-rank}{约化秩}\Term{locally-constant}{局部常值}（见 1.7 b)）。则 $G$ 存在\Term{reductive-center}{约化中心} $Z$。若 $G/Z$ \Term{representable}{可表示}（例如 $G$ 在 $S$ 上\Term{affine}{仿射}），则 $G$ 的\Term{maximal-torus}{极大环面} $T$（相应地，\Term{cartan-subgroup}{Cartan 子群} $C$）与 $G'=G/Z$ 的\Term{maximal-torus}{极大环面} $T'$（相应地，\Term{cartan-subgroup}{Cartan 子群} $C'$）一一对应；$T$（相应地，$C$）对应于 $T'=T/Z$（相应地，$C'=C/Z$），而 $T'$（相应地，$C'$）对应于 $T=\varphi^{-1}(T')$（相应地，$C=\varphi^{-1}(C')$），其中 $\varphi:G\to G'$ 为\Term{canonical-homomorphism}{典范同态}。

 d) 设 $T$ 是 $G$ 的\Term{maximal-torus}{极大环面}，以 $\mathfrak g$ 表示 $G$ 的\Term{lie-algebra}{Lie 代数}，考虑 $G$ 的\Term{adjoint-representation}{伴随表示}（II 4）所诱导的\Term{homomorphism}{同态}
\[\theta:T\longrightarrow\mathrm{GL}(\mathfrak g).\]
则 $\theta$ 的核是 $G$ 的\Term{reductive-center}{约化中心}。
\end{SGAresult}
\begin{proof}
a)、b)。a) 及 b) 第一个断言的证明与 1.7 a)、b) 的证明实质相同，这里不再重复论证。只指出，a) 需要用 IX 5.6（利用 $G$ 的\Term{fiber}{纤维}\Term{connected}{连通}这一事实）。证明 b) 的第二个断言：若 $Z$ 是 $G$ 的\Term{reductive-center}{约化中心}，则 $G'=G/Z$ 以单位\Term{subgroup}{子群}为\Term{reductive-center}{约化中心}。首先注意，由 VIII 5.1，\Term{quotient-group}{商群} $G/Z$ 确实\Term{representable}{可表示}，它在 $S$ 上\Term{affine}{仿射}、\Term{finite-presentation}{有限表示}（VIII 5.8），甚至在 $S$ 上\Term{smooth}{光滑}（因 $G'$ 在 $S$ 上\Term{flat}{平坦}且\Term{finite-presentation}{有限表示}，只须对\Term{fiber}{纤维}验证；而关于\Term{fiber}{纤维}，在 VI$_{\mathrm B}$ 9.2 (xii) 已指出，域 $k$ 上\Term{smooth}{光滑}\Term{algebraic-group}{代数群}的商仍\Term{smooth}{光滑}）。最后，由于 $G$ 的\Term{fiber}{纤维}\Term{connected}{连通}，$G'$ 也如此。因此 $G'$ 满足与 $G$ 相同的初始假设。要看出 $G'$ 的单位\Term{subgroup}{子群}是\Term{reductive-center}{约化中心}，若愿意可只考虑 $S$ 为域的谱的情形（4.3）。设 $Z'$ 为 $G'$ 的\Term{central}{中心}\Term{multiplicative-type}{乘法型}\Term{subgroup}{子群}，只须证明它就是\Term{unit-group}{单位群}。设 $Z_1$ 是 $Z'$ 在 $G$ 中的\Term{inverse-image}{逆像}，并考虑 $G$ 的\Term{inner-automorphism}{内自同构}所诱导的 $Z_1$ 在 $G$ 上的\Term{action}{作用}。因为 $Z$ 为\Term{central}{中心}\Term{subgroup}{子群}，其\Term{action}{作用}平凡，所以 $Z_1$ 的\Term{action}{作用}经由 $Z'$ 在 $G$ 上的\Term{action}{作用}给出。又因 $Z$ 在 $G$ 的\Term{center}{中心}内，$Z_1$ 从而 $Z'$ 在 $Z$ 上的\Term{action}{作用}平凡；此外因 $Z'$ 在 $G'$ 的\Term{center}{中心}内，$Z'$ 在商 $G/Z=G'$ 上的\Term{action}{作用}也平凡。由于 $Z$ 是\Term{central}{中心}的，立即\Term{pushout}{推出} $Z'$ 在 $G$ 上的\Term{action}{作用}来自\Term{group-homomorphism}{群同态}
\[u:Z'\longrightarrow\underline{\operatorname{Hom}}_{S\text{-gr}}(G',Z),\]
具体由
\[r(z')\cdot g=g\cdot u(z')(\varphi(g))\]
给出，其中 $r:Z'\to\underline{\operatorname{Aut}}_{S\text{-gr}}(G)$ 是所考虑的 $Z'$ 的 $G$-\Term{automorphism}{自同构}表示，$\varphi:G\to G'$ 为\Term{canonical-homomorphism}{典范同态}。
\SourcePage{157}
给定上面的\Term{group-homomorphism}{群同态} $u$，\Term{equivalence}{等价}于给定\Term{group-homomorphism}{群同态}
\[v:G'\longrightarrow\underline{\operatorname{Hom}}_{S\text{-gr}}(Z',Z).\]
另一方面，由 X 5.8，右端\Term{representable}{可表示}，且为 $S$ 上的\Term{etale}{平展}群，故其单位\Term{section}{截面}为\Term{open-immersion}{开浸入}；所以 $\operatorname{Ker}v$ 为 $G'$ 的\Term{open-subgroup}{开子群}。由于 $G'$ 的\Term{fiber}{纤维}\Term{connected}{连通}，它等于 $G'$，所以 $v$ 为零，故 $u$ 为零，故 $Z'$ 在 $G$ 上\Term{action}{作用}平凡；$Z_1$ 也如此，所以 $Z_1$ 是 $G$ 的\Term{central}{中心}\Term{subgroup}{子群}。因此 $Z_1$ 是\Term{multiplicative-type}{乘法型}群 $Z'$ 被\Term{multiplicative-type}{乘法型}群 $Z$ 的交换\Term{extension}{扩张}，故（因\Term{base}{基底}为域）仍为\Term{multiplicative-type}{乘法型}群（7.1.1）。因为它在 $G$ 的\Term{center}{中心}内，必包含在\Term{reductive-center}{约化中心} $Z$ 中，即 $Z_1=Z$，从而 $Z'$ 为\Term{unit-group}{单位群}。这完成了 b) 的证明。

 d) 设 $Z=\operatorname{Ker}\theta$，这是 $T$ 的一个\Term{multiplicative-type}{乘法型}\Term{subgroup}{子群}（例如由 IX 6.8）。由 4.5，任意\Term{central}{中心}\Term{homomorphism}{同态} $u:H\to G$，其中 $H$ 为\Term{multiplicative-type}{乘法型}且\Term{finite-type}{有限型}，都经由 $T$，从而经由 $Z$ \Term{resolution}{分解}。由于 $Z$ 的形成与任意\Term{base-change}{基变换}相容，只须证明 $Z$ 是\Term{central}{中心}的，即 $Z$ 的\Term{centralizer}{中心化子} $C$ 等于 $G$。由 XI 5.3，$C$ 是 $G$ 的\Term{smooth}{光滑}\Term{closed}{闭}\Term{subgroup}{子群}；另一方面，因为 $Z$ 在 $\mathfrak g$ 上的\Term{action}{作用}平凡，有 $\operatorname{Lie}(C)=\mathfrak g$。容易由此\Term{pushout}{推出} $C=G$：事实上，\Term{immersion}{浸入} $C\to G$ 是\Term{etale}{平展}的，因为它逐\Term{fiber}{纤维}\Term{etale}{平展}（作为同维\Term{smooth}{光滑}\Term{algebraic-group}{代数群}之间的\Term{unramified}{非分歧}\Term{homomorphism}{同态}，VI$_{\mathrm B}$ 1.3），然后可以应用 X 3.5。因此 $C\to G$ 是\Term{etale}{平展}\Term{closed-immersion}{闭浸入}，因而是\Term{open-immersion}{开浸入}（SGA 1 I 5.1）；它同时也是\Term{closed-immersion}{闭浸入}，而 $G$ 的\Term{fiber}{纤维}\Term{connected}{连通}，故它是\Term{isomorphism}{同构}。

 c) 由 1.7 b)，$G$ 按\Term{fpqc-topology}{fpqc 拓扑}\Term{local-scheme}{局部}地存在\Term{maximal-torus}{极大环面}；故由 4.2 及刚证明的 d)，$G$ 存在\Term{reductive-center}{约化中心} $Z$。已在 d) 中看到，$G$ 的每个\Term{maximal-torus}{极大环面} $T$ 都包含 $Z$。立即可见 $T'=T/Z$ 是\Term{torus}{环面}；要证明它是 $G'$ 中的\Term{maximal-torus}{极大环面}，由定义 1.3 可归结到 $S$ 为\Term{algebraically-closed-field}{代数闭域}的谱，此时断言包含在 BIBLE 7 \Term{theorem}{定理} 3 a) 中。反过来，设 $T'$ 为 $G'$ 的\Term{maximal-torus}{极大环面}，令 $T=\varphi^{-1}(T')$，要证明 $T$ 是 $G$ 的\Term{maximal-torus}{极大环面}。问题对于\Term{fpqc-topology}{fpqc 拓扑}是\Term{local-scheme}{局部}的，故可设 $G$ 已有一个\Term{maximal-torus}{极大环面} $T_0$，于是由前述结论，$T'_0=T_0/Z_0$ 是 $G'$ 的\Term{maximal-torus}{极大环面}。由\Term{conjugacy}{共轭}\Term{theorem}{定理} 1.7 c)，$T'$、$T'_0$ 按\Term{fpqc-topology}{fpqc 拓扑}\Term{local-scheme}{局部}\Term{conjugacy}{共轭}；所以（由于 $\varphi:G\to G'$ 对此\Term{topology}{拓扑}是\Term{covering}{覆盖}，故 $G'$ 的每个\Term{section}{截面}\Term{local-scheme}{局部}地\Term{liftable}{可提升}为 $G$ 的\Term{section}{截面}）$T$ 和 $T_0$ 也\Term{local-scheme}{局部}\Term{conjugacy}{共轭}。因为 $T$ 是\Term{maximal-torus}{极大环面}，故 $T$ 也如此。对于\Term{cartan-subgroup}{Cartan 子群}可类似处理。证明完毕。
\end{proof}
当 $T$ \Term{diagonalizable}{可对角化}，因而形如 $\mathrm D_S(M)$，其中 $M$ 为\Term{finite-type}{有限生成}自由 $\mathbf Z$-模时，给出 d) 的一个有用表述。此时（I 4.7.3），$T$-模
\SourcePage{158}
$\mathfrak g$ \Term{resolution}{分解}为子 $T$-模 $\mathfrak g_m$（$m\in M$）的\Term{direct-sum}{直和}：
\[\mathfrak g=\coprod_{m\in M}\mathfrak g_m,\]
它们必为\Term{locally-free}{局部自由}。假定对每个 $m\in M$，$\mathfrak g_m$ 的秩为常数（特别地，当 $S$ \Term{connected}{连通}时如此）。于是满足 $\mathfrak g_m\ne0$ 的 $m\in M$ 所成\Term{set}{集合} $R$（“\Term{root}{根}”）\Term{finite}{有限}。于是：
\begin{SGAresult}{推论}{4.8}
在上述条件及记号下，$G$ 的\Term{reductive-center}{约化中心}是 $T$ 上各\Term{root-character}{根特征} $m\in R$ 的核的交。因此有\Term{isomorphism}{同构}
\[Z\simeq\mathrm D_S(N),\]
其中 $N$ 为 $M$ 除以由 $R$ 生成的\Term{subgroup}{子群}所得的商。
\end{SGAresult}
\begin{SGAresult}{推论}{4.9}
设 $G$ 是\Term{algebraically-closed-field}{代数闭域} $k$ 上的\Term{algebraic-group}{代数群}。假定 $G$ 在 $k$ 上\Term{smooth}{光滑}，\Term{connected}{连通}、\Term{affine}{仿射}，\Term{reductive-center}{约化中心}为\Term{unit-group}{单位群}，且 $G$ 的\Term{lie-algebra}{Lie 代数} $\mathfrak g$ \Term{nilpotent}{幂零}。则 $G$“\Term{unipotent}{幺幂}”，即有一个商\Term{isomorphism}{同构}于 $\mathbf G_a$ 的\Term{composition-series}{合成列}。
\end{SGAresult}
\begin{proof}
由 BIBLE 6 \Term{theorem}{定理} 4 推论 3，只须证明 $G$ 的\Term{maximal-torus}{极大环面} $T$ 为\Term{unit-group}{单位群}，\Term{equivalence}{等价}地，$T$ 的\Term{lie-algebra}{Lie 代数} $\mathfrak t$ 为 $0$。由于 $T$-模 $\mathfrak g$ 按 $T$ 的\Term{character}{特征} $\alpha$ \Term{resolution}{分解}（I 4.7.3），对每个 $t\in\mathfrak t$，$\mathfrak g$ 上的\Term{operator}{算子} $\operatorname{ad}_{\mathfrak g}(t)$ 是\Term{semisimple}{半单}的。因此若 $\mathfrak g$ \Term{nilpotent}{幂零}，则 $\operatorname{ad}_{\mathfrak g}(t)$ 为零。又由 4.7 d)，因为 $G$ 的\Term{reductive-center}{约化中心}只有单位元，\Term{homomorphism}{同态} $T\to\mathrm{GL}(\mathfrak g)$ 为\Term{monomorphism}{单态射}，从而在\Term{lie-algebra}{Lie 代数}上诱导\Term{injectivity}{单射}，即（II 4.5）对 $t\in\mathfrak t$，关系 $\operatorname{ad}_{\mathfrak g}(t)=0$ 蕴含 $t=0$。这证明 $\mathfrak t=0$，证明完毕。
\end{proof}
\begin{SGAresult}{命题}{4.10}
设 $G$ 是\Term{algebraically-closed-field}{代数闭域} $k$ 上\Term{affine}{仿射}、\Term{smooth}{光滑}、\Term{connected}{连通}的\Term{algebraic-group}{代数群}。则 $G$ 的\Term{reductive-center}{约化中心}是 $G$ 的所有\Term{maximal-torus}{极大环面}之交。
\end{SGAresult}
当然，这里的交是\Term{scheme}{概形}意义下的交（\Term{equivalence}{等价}地，作为 $G$ 的\Term{subfunctor}{子函子}的交），即包含于 $G$ 的每个\Term{maximal-torus}{极大环面}中的最大\Term{closed}{闭}\Term{subprescheme}{子预概形}。由 $G$ 为 Noether 的，它也是适当选取的\Term{finite}{有限}多个\Term{maximal-torus}{极大环面}的交。
\begin{proof}
设 $Z$ 为所述交。$Z$ 是某\Term{torus}{环面}的\Term{closed}{闭}\Term{subgroup}{子群}，所以为\Term{multiplicative-type}{乘法型}；另一方面，由 4.5，它包含 $G$ 的\Term{reductive-center}{约化中心}。要证明二者相等，只须证明 $Z$ 是\Term{central}{中心}的。由于 $G$ \Term{connected}{连通}，由 IX 5.5，只须证明 $Z$ 为\Term{normal-subgroup}{正规}\Term{subgroup}{子群}。由构造，$Z$ 在所有 $g\in G(k)$ 对应的 $\operatorname{int}(g)$ 下\Term{invariant}{不变}，故 $Z$ 的\Term{normalizer}{正规化子} $N$（见 VIII 6.7）是包含 $G$ 的所有\Term{rational-point}{有理点}的\Term{closed}{闭}\Term{subgroup}{子群}。因为 $G$ \Term{reduced}{既约}，有 $N=G$，证明完毕。
\end{proof}
\begin{SGAresult}{命题}{4.11}
设 $S$ 是\Term{prescheme}{预概形}，$G$ 是 $S$ 上\Term{smooth}{光滑}、\Term{affine}{仿射}、\Term{fiber}{纤维}\Term{connected}{连通}且\Term{unipotent-rank}{幺幂秩}为零的 $S$-\Term{group-prescheme}{群预概形}。由 1.9，
\SourcePage{159}
$G$ 的\Term{reductive-rank}{约化秩}\Term{locally-constant}{局部常值}，所以由 1.10，$G$ 的“\Term{maximal-torus-prescheme}{极大环面预概形}” $\mathcal T$ 有定义，且在 $S$ 上\Term{smooth}{光滑}、\Term{separated}{分离}、\Term{finite-type}{有限型}。让 $G$ 通过\Term{inner-automorphism}{内自同构}\Term{action}{作用}在 $\mathcal T$ 上，得到\Term{group-functor}{群函子}\Term{homomorphism}{同态}
\[u:G\longrightarrow\underline{\operatorname{Aut}}_S(\mathcal T).\]
在这些条件下，$G$ 的以下三个\Term{subfunctor}{子函子}相同：

 (i) $G$ 的\Term{reductive-center}{约化中心} $Z$。

 (ii) $G$ 的\Term{center}{中心} $Z'$。

 (iii) 上述\Term{homomorphism}{同态}的核 $Z''=\operatorname{Ker}u$。

特别地，$G$ 的\Term{center}{中心}由 $G$ 的\Term{multiplicative-type}{乘法型}\Term{subgroup}{子群}\Term{representable}{表示}。
\end{SGAresult}
\begin{proof}
显然 $Z\subset Z'\subset Z''$，只须证明 $Z''\subset Z$。由于假设对\Term{base-change}{基变换}\Term{stable}{稳定}，这\Term{equivalence}{等价}于证明：$G$ 在 $S$ 上任何在 $\mathcal T$ 上\Term{action}{作用}平凡的\Term{section}{截面} $g$ 都是 $Z$ 的\Term{section}{截面}。引入 $G'=G/Z$ 并使用 4.7 b)、c)（它们特别蕴含 $G'$ 的\Term{maximal-torus}{极大环面}\Term{prescheme}{预概形} $\mathcal T'$ 与 $\mathcal T$ \Term{canonical-isomorphism}{典范同构}），可归结到 $G=G'$，即 $G$ 的\Term{reductive-center}{约化中心} $Z$ 为零的情形。（注意，由 4.7 c)，$G'$ 的\Term{unipotent-rank}{幺幂秩}等于 $G$ 的\Term{unipotent-rank}{幺幂秩}，故为零。）此时要证明 $g$ 为 $G$ 的单位\Term{section}{截面}。通常的化简方法将问题归结到 $S$ 为 Noether 的，甚至 $S$ 为 Artin \Term{local-scheme}{局部}的情形（因为当 $S$ \Term{locally-noetherian}{局部 Noether} 时，要验证 $g$ 为单位\Term{section}{截面}，只须在每个 $S'$ 为 Artin \Term{local-scheme}{局部}的\Term{base-change}{基变换} $S'\to S$ 后验证）。此时 $u$ 的核 $Z''$ \Term{representable}{可表示}（VIII 6.2 a) 及 6.5 c)）。（注意，由 XI 6.8，$Z''$ \Term{representable}{可表示}。）由 \Term{nakayama-lemma}{Nakayama 引理}，要证明 $Z''$ 为单位\Term{subgroup}{子群}，只须证明它的\Term{fiber}{纤维} $Z''_0$ 如此。这又归结到 $S$ 为域 $k$ 的谱的情形，显然可设 $k$ \Term{algebraically-closed}{代数闭}。$Z''$ 包含于 $\mathcal T(k)$ 每个点的\Term{stabilizer}{稳定子}中，即包含于 $G$ 每个\Term{maximal-torus}{极大环面} $T$ 的\Term{normalizer}{正规化子} $N(T)$ 中。因为 $G$ 的\Term{reductive-rank}{约化秩}为零，由 XI 5.9，$T$ 是 $N(T)$ 的\Term{open-subgroup}{开子群}，故 $N(T)$ 的\Term{lie-algebra}{Lie 代数}与 $T$ 的相同。因此 $Z''$ 的\Term{lie-algebra}{Lie 代数}包含于 $T$ 的\Term{lie-algebra}{Lie 代数}。另一方面，由 4.10，各\Term{maximal-torus}{极大环面} $T$ 的\Term{lie-algebra}{Lie 代数}之交正是\Term{reductive-center}{约化中心} $Z$ 的\Term{lie-algebra}{Lie 代数}，这里因为假定 $Z=0$，它为零。因此 $Z''$ 的\Term{lie-algebra}{Lie 代数}为零，即 $Z''$ 在 $k$ 上\Term{etale}{平展}。又 $Z''$ 显然在 $G$ 中\Term{normal-subgroup}{正规}，由于 $G$ \Term{connected}{连通}，容易\Term{pushout}{推出} $Z''$ 在 $G$ 的\Term{center}{中心}内。因此对每个\Term{maximal-torus}{极大环面} $T$，它都包含在 $T=\operatorname{Centr}_G(T)$ 中，从而包含在所有\Term{maximal-torus}{极大环面}之交中，即由 4.10 包含在 $Z=0$ 中。证明完毕。
\end{proof}
\SGAsection{5}{对乘法型子群概形的应用}
\footnote{原注（*）：本节结果在不假定 $G$ 在 $S$ 上\Term{affine}{仿射}时的推广，见 M. Raynaud：《\Term{group-scheme}{群概形}和\Term{homogeneous-space}{齐性空间}上的\Term{ample-sheaf}{丰沛层}》（Faisceaux amples sur les schémas en groupes et les espaces homogènes），第 IX 章。}
\begin{SGAresult}{\Term{theorem}{定理}}{5.1}
设 $S$ 是\Term{prescheme}{预概形}，$G$ 是 $S$ 上\Term{smooth}{光滑}且\Term{affine}{仿射}的 $S$-\Term{group-prescheme}{群预概形}，$M$ 是“$G$ 的\Term{multiplicative-type}{乘法型}\Term{subgroup}{子群}\Term{prescheme}{预概形}”，表示 XI 4.1 中明确给出的\Term{functor}{函子}。
\SourcePage{160}
对每个整数 $n>0$，令 $T_n$ 为 $M$ 的\Term{subfunctor}{子函子}，其 $T_n(S')$ 等于满足 $n\operatorname{id}_H=0$ 的 $G_{S'}$ 的\Term{multiplicative-type}{乘法型}\Term{group-subprescheme}{群子预概形} $H$ 的\Term{set}{集合}；于是 $T_n$ \Term{representable}{可表示}，且在 $S$ 上\Term{affine}{仿射}（XI 3.12 a)）。设 $u_n:M\to T_n$ 为 $u_n(H)={}_nH$ 所定义的\Term{morphism}{态射}（其中 ${}_nH=\operatorname{Ker}(n\operatorname{id}_H)$）。在这些记号下：

 a) $M$ 的每个在 $S$ 上\Term{finite-type}{有限型}的\Term{subprescheme}{子预概形} $U$ 都包含在一个在 $S$ 上\Term{finite-type}{有限型}的\Term{closed}{闭}\Term{subprescheme}{子预概形}中；且 $M$ 的每个在 $S$ 上\Term{finite-type}{有限型}的\Term{closed}{闭}\Term{subprescheme}{子预概形} $X$ 都在 $S$ 上\Term{affine}{仿射}。

 b) 假定 $S$ \Term{quasi-compact}{拟紧}，且 $X$ 是 $M$ 的一个在 $S$ 上\Term{finite-type}{有限型}的\Term{closed}{闭}\Term{subprescheme}{子预概形}。则存在整数 $n>0$，使对 $n$ 的任意倍数 $m$，诱\Term{derived-morphism}{导态射}
\[u_m|_X:X\longrightarrow T_m\]
是\Term{closed-immersion}{闭浸入}。
\end{SGAresult}
\begin{proof}
a) 要证明 a) 的第一个断言，取 $U$ 的\Term{schematic-closure}{概形闭包}为 $X$（EGA I 9.5.1 和 9.5.3）；由于\Term{immersion}{浸入} $i:U\to M$ 是\Term{quasi-compact}{拟紧}的（因 $U$ 在 $S$ 上\Term{finite-type}{有限型}且 $M$ 在 $S$ 上\Term{separated}{分离}，见 4.1），此\Term{closure}{闭包}有定义。所以须证明这样的 $X$ 在 $S$ 上\Term{affine}{仿射}，这也将同时证明 a) 的第二个断言。由此表述，可见问题在 $S$ 上是\Term{local-scheme}{局部}的，因此可设 $S$ \Term{affine}{仿射}。此时 $U$ 在 $S$ 上\Term{finite-type}{有限型}，故包含在 $M$ 的某个\Term{quasi-compact}{拟紧}\Term{open-subset}{开集}中；从而可归结到 $U$ 自身为在 $S$ 上\Term{finite-type}{有限型}的\Term{open-subset}{开集}，即\Term{quasi-compact}{拟紧}\Term{open-subset}{开集}的情形。（注意：$S$ 上\Term{affine}{仿射}\Term{scheme}{概形}的\Term{closed}{闭}\Term{subscheme}{子概形}仍在 $S$ 上\Term{affine}{仿射}。）

只须证明这样的 $U$ 包含在 $M$ 的某个\Term{affine}{仿射}\Term{closed}{闭}\Term{subprescheme}{子预概形}中。按这种\Term{form}{形式}，通常的化简方法立即将问题归结到 $S$ 为 Noether 的情形。对 b) 同样处理，可归结到 $S$ \Term{affine}{仿射}且 Noether 的情形。

重新考虑 XI 4.1 证明中使用的 $T_n$ 的\Term{projective-limit}{投射极限} $T$。它是 $S$ 上的\Term{affine}{仿射}\Term{prescheme}{预概形}（但并非\Term{finite-type}{有限型}），而且其\Term{local-ring}{局部环}都是 Noether 环，正如 IX 3.7 证明开头所见。（注意：对 $t\in T$，$t=(t_n)$，因为\Term{transition-morphism}{过渡态射} $T_m\to T_n$ \Term{smooth}{光滑}，且 $T_m$ 的\Term{dimension}{维数}有上界，存在 $n_0$，使对 $n_0$ 的每个倍数 $n$，$T_n\to T_{n_0}$ 在 $t_n$ 处\Term{etale}{平展}。）该处的证明或 XI 3.11 表明，\Term{canonical}{典范}\Term{morphism}{态射}
\[u:M\longrightarrow T\]
是\Term{immersion}{浸入}，且在\Term{local-ring}{局部环}上\Term{induction-isomorphism}{诱导同构}（但应注意，$u$ 一般既非\Term{open-immersion}{开浸入}，也非\Term{quasi-compact}{拟紧}\Term{morphism}{态射}）。设 $U$ 是 $M$ 的\Term{quasi-compact}{拟紧}\Term{open-subset}{开集}，我们将证明它在 $T$ 中的\Term{closure}{闭包}包含在 $M$ 中。这说明（若 $X$ 表示 $U$ 在 $M$ 中的\Term{schematic-closure}{概形闭包}）$u|_X:X\to T$ 是\Term{closed-immersion}{闭浸入}，所以 $X$
\SourcePage{161}
\Term{affine}{仿射}，从而证明 a)。由于 $U$ 为 Noether 的，它只有\Term{finite}{有限}多个\Term{irreducible-component}{不可约分支}，且 $U$ 的\Term{closure}{闭包}中的每个元素 $t\in T$ 都是 $U$ 中某个元素 $x$ 的\Term{specialization}{特殊化}。因为 $t$ 在 $T$ 中的\Term{local-ring}{局部环}为 Noether 环，$t$ 在赋予诱导\Term{reduced}{既约}结构的 $\overline{\{x\}}$ 中的\Term{local-ring}{局部环} $A$ 也为 Noether 环。

由整 Noether \Term{local-scheme}{局部}\Term{scheme}{概形} $S'=\operatorname{Spec}(A)$ 到 $T$ 的\Term{canonical}{典范}\Term{morphism}{态射}，将 $S'$ 的\Term{closed-point}{闭点}送到 $t$，将\Term{generic-point}{一般点}送到 $x\in M$。须在这些条件下证明 $t\in M$。以适当的在 $A$ 上\Term{flat}{平坦}的\Term{complete}{完备}\Term{local-scheme}{局部} $A$-\Term{algebra}{代数}（EGA $0_{\mathrm{III}}$ 10.3.1）除以一个\Term{minimal-prime-ideal}{极小素理想}所得的商替换 $A$ 后，可设 $A$ \Term{complete}{完备}且\Term{residue-field}{剩余域}\Term{algebraically-closed}{代数闭}（借助 EGA II 7.1.7，甚至还可归结到它为\Term{discrete-valuation-ring}{离散赋值环}的情形）。于是（改记号，将 $S'$ 记为 $S$）归结到：
\end{proof}
\begin{SGAresult}{\Term{lemma}{引理}}{5.2}
设 $S$ 是\Term{residue-field}{剩余域}\Term{algebraically-closed}{代数闭}的\Term{complete}{完备}整 \Term{noetherian-local-ring}{Noether 局部环}的谱，$G$ 是 $S$ 上\Term{smooth}{光滑}且\Term{affine}{仿射}的 $S$-\Term{group-prescheme}{群预概形}，$(H(n))_{n>0}$ 是 $G$ 的一个\Term{multiplicative-type}{乘法型}\Term{subgroup}{子群}系统，以整数 $n>0$ 编号，$\eta$ 是 $S$ 的\Term{generic-point}{一般点}。假定：

 a) 若 $m$ 为 $n$ 的倍数，则 $H(n)={}_nH(m)$。

 b) 存在 $G_\eta$ 的一个\Term{multiplicative-type}{乘法型}\Term{subgroup}{子群} $H_\eta$，使对每个 $n>0$，有 $H(n)_\eta={}_n(H_\eta)$。

在这些条件下，存在 $G$ 的一个\Term{multiplicative-type}{乘法型}\Term{subgroup}{子群} $H$，使对每个 $n>0$，有 $H(n)={}_nH$。
\end{SGAresult}
\begin{proof}
对每个 $n$，令 $C_n=\operatorname{Centr}_G(H(n))$，这是 $G$ 的一个\Term{closed}{闭}\Term{group-subprescheme}{群子预概形}，在 $S$ 上\Term{smooth}{光滑}（XI 5.3）。将整数 $n>0$ 按整除关系排序，则 $C_n$ 构成 $G$ 的\Term{closed}{闭}\Term{subprescheme}{子预概形}的递减族；因为 $G$ 为 Noether 的，此族最终\Term{stable}{稳定}。令 $C$ 为充分大 $n$ 时 $C_n$ 的公共值。则 $C$ 满足与 $G$ 相同的条件，且 $H(n)$ 事实上都是 $C$ 的\Term{subgroup}{子群}。所以以 $C$ 替换 $G$ 后，可设 $H(n)$ 是 $G$ 的\Term{central}{中心}\Term{subgroup}{子群}。

设 $s$ 是 $S$ 的\Term{closed-point}{闭点}，$Z_s$ 是 $G_s$ 的\Term{reductive-center}{约化中心}（由 4.4 定义）。由定义 4.1，$Z_s$ 包含 $H(n)_s$。由于 $A$ \Term{complete}{完备}，$Z_s$ 来自 $G$ 的某个\Term{multiplicative-type}{乘法型}\Term{group-subprescheme}{群子预概形} $Z$（XI 5.8）。我断言 $Z$ 包含各 $H(n)$。事实上，由于 $H(n)$ 是\Term{central}{中心}的，所以它在 $G$ 中与 $Z$ 交换，得到\Term{group-homomorphism}{群同态} $Z\times H(n)\to G$。由 IX 6.8，它有像群，即 $G$ 的一个包含 $Z$ 的\Term{multiplicative-type}{乘法型}\Term{subgroup}{子群} $K$；只须证明 $K=Z$，而这由 $K_s=Z_s$ 及 IX 5.1 bis \Term{pushout}{推出}。

于是只须证明 5.2 的对应命题，其中将 $G$ 换成在 $S$ 上\Term{multiplicative-type}{乘法型}且\Term{finite-type}{有限型}的群 $Z$（不必在 $S$ 上\Term{smooth}{光滑}）。由于 $A$ \Term{complete}{完备}且\Term{residue-field}{剩余域}\Term{algebraically-closed}{代数闭}，$Z$ \Term{diagonalizable}{可对角化}（X 3.3 和 1.4），所以形如 $\mathrm D_S(M)$，其中 $M$ 为\Term{finite-type}{有限生成}\Term{abelian-group}{交换群}。因此 $Z$ 的任意\Term{multiplicative-type}{乘法型}\Term{subgroup}{子群} $H$ 都\Term{diagonalizable}{可对角化}（IX 2.11 (i)），所以形如 $\mathrm D_S(N)$，其中 $N$ 是 $M$ 的\Term{quotient-group}{商群}（VIII 3.2）。于是各 $H(n)$ 对应于 $M$ 的\Term{quotient-group}{商群} $M(n)$，而存在 $Z$ 的一个\Term{multiplicative-type}{乘法型}\Term{subgroup}{子群}
\SourcePage{162}
$H$，使对每个 $n$，有 $H(n)={}_nH$，就\Term{equivalence}{等价}于存在 $M$ 的\Term{quotient-group}{商群} $N$，使对每个 $n$，有 $M(n)=N/nN$。这立即由存在 $G_\eta$ 的\Term{subgroup}{子群} $H_\eta$、使对每个 $n$ 有 $H(n)_\eta={}_n(H_\eta)$ 这一事实\Term{pushout}{推出}。由此完成了 5.2 从而 a) 的证明。
\end{proof}
\begin{proof}[\Term{theorem}{定理} 5.1 的证明（续）]
b) 已知 $u|_X:X\to T$ 为\Term{closed-immersion}{闭浸入}。由此容易\Term{pushout}{推出}，当 $m$ 充分大时，\Term{composition-morphism}{复合态射} $u_m|_X:X\to T_m$ 也是\Term{closed-immersion}{闭浸入}。以后不需要此事实，故这里不详细展开证明。
\end{proof}
\begin{SGAresult}{推论}{5.3}
沿用 5.1 的记号，设 $U$ 为 $M$ 中既开又闭、在 $S$ 上\Term{finite-type}{有限型}的部分。则 $U$ 赋予 $M$ 诱导的结构后在 $S$ 上\Term{affine}{仿射}；若 $S$ \Term{quasi-compact}{拟紧}，则存在整数 $n>0$，使对 $n$ 的每个倍数 $m$，诱\Term{derived-morphism}{导态射} $u_m|_U:U\to T_m$ 是\Term{open-and-closed}{开闭}\Term{immersion}{浸入}（即到 $T_m$ 的一个赋予\Term{induced-structure}{诱导结构}的既开又闭部分的\Term{isomorphism}{同构}）。
\end{SGAresult}
第一个断言由 5.1 a) \Term{pushout}{推出}；第二个由 5.1 b) \Term{pushout}{推出}，只须注意 $u_m:M\to T_m$ \Term{smooth}{光滑}（XI 2.2 bis），而\Term{smooth}{光滑}即\Term{etale}{平展}的\Term{immersion}{浸入}是\Term{open-immersion}{开浸入}。
\begin{SGAresult}{推论}{5.4}
设 $S$ 是\Term{prescheme}{预概形}，$G$ 是 $S$ 上\Term{smooth}{光滑}、\Term{affine}{仿射}且\Term{reductive-rank}{约化秩}\Term{locally-constant}{局部常值}（1.7 b)）的 $S$-\Term{group-prescheme}{群预概形}，$\mathcal T$ 为“$G$ 的\Term{maximal-torus-prescheme}{极大环面预概形}”（1.10）。则 $\mathcal T$ 在 $S$ 上\Term{smooth}{光滑}且\Term{affine}{仿射}。若 $T$ 是 $G$ 的\Term{maximal-torus}{极大环面}，$N(T)$ 是其\Term{normalizer}{正规化子}，则 $G/N(T)$（XI 5.3 bis）在 $S$ 上\Term{affine}{仿射}。若 $W(T)=N(T)/C(T)$ 在 $S$ 上\Term{finite}{有限}（见 2.1 b)），则 $G/C(T)$ 也如此，其中 $C(T)$ 为 $T$ 的\Term{centralizer}{中心化子}。
\end{SGAresult}
\begin{proof}
第二个断言包含在第一个中，因为由\Term{conjugacy}{共轭}\Term{theorem}{定理}，$G/N(T)$ \Term{isomorphism}{同构}于 $\mathcal T$（XI 5.5 bis）。对于第一个断言，注意由构造，$\mathcal T$ \Term{isomorphism}{同构}于 $M$ 的一个既开又闭的\Term{subprescheme}{子预概形}（因为可设 $G$ 的\Term{reductive-rank}{约化秩}为常数 $r$，于是 $\mathcal T$ 是 $M$ 中对应于\Term{relative-dimension}{相对维数}为 $r$ 的\Term{subtorus}{子环面}的\Term{subprescheme}{子预概形}，即 $M$ 中使“泛\Term{multiplicative-type}{乘法型}\Term{subgroup}{子群}” $H\subset G_M$ 为\Term{relative-dimension}{相对维数} $r$ 的\Term{torus}{环面}的最大\Term{subprescheme}{子预概形}）。因此可应用 5.3。最后，对最后一个断言，只须注意 $G/C(T)$ 在 $G/N(T)\simeq(G/C(T))/W(T)$ 上\Term{finite}{有限}，故因 $G/N(T)$ \Term{affine}{仿射}而\Term{affine}{仿射}。
\end{proof}
\begin{SGAresult}{推论}{5.5}
设 $G$ 是域 $k$ 上\Term{smooth}{光滑}且\Term{affine}{仿射}的\Term{algebraic-group}{代数群}。则 $G$ 的\Term{multiplicative-type}{乘法型}\Term{subgroup}{子群}\Term{scheme}{概形} $M$ 是 $S$ 上\Term{affine}{仿射}\Term{scheme}{概形}的\Term{direct-sum}{直和}。对 $G$ 的每个\Term{multiplicative-type}{乘法型}\Term{subgroup}{子群} $H$，若 $C$、$N$ 分别为它的\Term{centralizer}{中心化子}、\Term{normalizer}{正规化子}，则商 $G/C$、$G/N$ 都\Term{affine}{仿射}。
\end{SGAresult}
\begin{proof}
由 XI 5.1 bis 可见，$M$ 的任何\Term{finite}{有限}\Term{closed-subset}{闭子集}在 $G$ 对 $M$ 的\Term{action}{作用}下的\Term{saturated}{饱和}集都是开的：事实上，可归结到 $k$ \Term{algebraically-closed}{代数闭}，从而归结到一个 $k$-\Term{rational-point}{有理点} $x$ 的\Term{orbit}{轨道}；此时由所引之处，$G\to M$，$g\mapsto g\cdot x$，为\Term{smooth}{光滑}\Term{morphism}{态射}，故为\Term{open-morphism}{开态射}，所以其像为\Term{open-subset}{开集}。
\SourcePage{163}
令 $U$ 为 $M$ 的各\Term{closed-point}{闭点}在 $G$ 的\Term{action}{作用}所定义的\Term{equivalence-relation}{等价关系}下的\Term{equivalence-class}{等价类}之并。则 $U$ 是\Term{open-subset}{开集}，且包含 $M$ 的所有\Term{closed-point}{闭点}，故由 Hilbert \Term{nullstellensatz}{零点定理}，$U=M$。因此 $M$ 是一些\Term{open-subset}{开集}的\Term{disjoint-union}{不交并}，这些\Term{open-subset}{开集}必然也是闭集，故 $M$ 是一些\Term{prescheme}{预概形} $M_i$ 的和\Term{prescheme}{预概形}；每个 $M_i$ 都是某个\Term{closed-point}{闭点}在 $G$ 下的\Term{orbit}{轨道}，所以\Term{quasi-compact}{拟紧}，从而\Term{finite-type}{有限型}。由 5.3，各 $M_i$ 都\Term{affine}{仿射}。若 $H$ 是 $G$ 的\Term{multiplicative-type}{乘法型}\Term{subgroup}{子群}，它对应于 $M$ 的一个 $k$-\Term{rational-point}{有理点} $x$，而 $G/N$ 等同于 $x$ 在 $G$ 下的\Term{orbit}{轨道}（XI 5.5 bis）。故由上述结论，$G/N$ \Term{affine}{仿射}。因为 $C$ 是 $N$ 的\Term{open-subgroup}{开子群}（XI 5.9），所以 $G/C$ 在 $G/N$ 上\Term{finite}{有限}（因为后者\Term{isomorphism}{同构}于 $(G/C)/(N/C)$），从而因 $G/C$ \Term{affine}{仿射}而\Term{affine}{仿射}。这个证明同时说明：
\end{proof}
\begin{SGAresult}{推论}{5.6}
对 $G$、$H$ 采用 5.5 的条件，则 $M$ 中“$G$ 的与 $H$ \Term{local-scheme}{局部}\Term{conjugacy}{共轭}的\Term{multiplicative-type}{乘法型}\Term{subgroup}{子群}”的\Term{subscheme}{子概形} $U$，是 $G$ 的一个既开又闭的\Term{subprescheme}{子预概形}。
\end{SGAresult}
形象地说，$G$ 的任何\Term{multiplicative-type}{乘法型}\Term{subgroup}{子群} $H'$，若是与 $H$ \Term{conjugacy}{共轭}的 $G$ 的\Term{subgroup}{子群}之\Term{limit}{极限}，则它自身也与 $H$ \Term{conjugacy}{共轭}。

\textbf{注记 5.7。}设 $S$ 是\Term{prescheme}{预概形}，$G$ 是 $S$ 上\Term{smooth}{光滑}且\Term{affine}{仿射}的 $S$-\Term{prescheme}{预概形}，$H$ 是\Term{multiplicative-type}{乘法型}且\Term{finite-type}{有限型}的 $S$-\Term{prescheme}{预概形}，置 $M=\underline{\operatorname{Hom}}_{S\text{-gr}}(H,G)$。由 XI 4.2，它\Term{representable}{可表示}，且在 $S$ 上\Term{smooth}{光滑}并\Term{separated}{分离}。可以对 $M$ 证明一个与 5.1 完全类似的结果，或以类似 XI 4.2 的证明归结到 5.1，或直接仿照 5.1 的证明。由此\Term{pushout}{推出} 5.3、5.5、5.6 的相应变体，留待读者表述。
